\documentclass[twoside,11pt]{article}

% Use JMLR style with preprint option
\usepackage[preprint]{jmlr2e}
\usepackage{amsmath}
\usepackage{bbm}
% Additional packages (jmlr2e already includes: epsfig, amssymb, natbib, graphicx)
\usepackage{subcaption}
\usepackage{xcolor}
\usepackage{tikz}
\usetikzlibrary{calc, arrows.meta, decorations.markings, patterns, shapes.geometric, positioning, 3d, decorations.pathmorphing}
\usepackage{array}
\usepackage{tabularx}
\usepackage{listings}
\lstset{basicstyle=\ttfamily, breaklines=true}
\usepackage{booktabs}
\usepackage{xr-hyper}

% Figure search paths (match main paper)
\graphicspath{{./}{./figs/}{../}{attention/}{../../Transformer Manuscript/fig_transformer_validation/}}

% Custom commands (must match main paper)
\newcommand{\softmax}{\operatorname{softmax}}
\newcommand{\KL}{\operatorname{KL}}
\newcommand{\Tr}{\operatorname{Tr}}
\externaldocument[][nocite]{GL(K)_attention_v2}

\setlength{\emergencystretch}{3em}
\begin{document}
\jmlrheading{}{2026}{}{}{}{}{Robert C. Dennis}

\title{Variational Belief Gradients for Gauge-Covariant Gaussian Attention\\[6pt]
\large Optional Objectives, Dynamics, and Numerical Methods}

\author{\name Robert C. Dennis \email cdenn016@gmail.com \\
  \addr Independent Researcher\\
  \addr Leander, Texas 78641, USA}


\maketitle

\ShortHeadings{Variational Belief Gradients}{Dennis}

This companion develops the differential theory of optional Gaussian belief revision. It specifies prior coupling, entropy-regularized and entropy-suppressed objectives, derivative support, covariance and frame gradients, and numerical update maps. The conditional forward-KL characterization concerns stationary belief updates under its stated assumptions. These constructions are separate from the projected routing and direct predictive-content operator developed in the main paper, \emph{Gauge-Covariant Gaussian Attention through Projected Matrix Links}. That predictor is trained by outer differentiation without requiring an inner gradient-based belief update.

The full-state alignment calculations below use the flat reference identification \(\Omega_{ij}=U_iU_j^{-1}\). Applying them to lower-dimensional comparison laws or general matrix links requires differentiating the corresponding declared objective; a change of notation alone does not establish that extension. The main paper's appendices retain the normalized-joint discussion and statistical conditions needed by the predictive architecture.

\renewcommand{\theequation}{G\arabic{equation}}
\renewcommand{\thefigure}{G\arabic{figure}}
\renewcommand{\thetable}{G\arabic{table}}

\section{Optional Belief-Revision Objectives and Derivative Support}
\label{app:optional_belief_revision}

The categorical routing coordinate and the update of a Gaussian belief are distinct mathematical operations. The main paper derives routing from projected comparison distributions and specifies an independent predictive content operator. This section studies an additional full-state belief-refinement objective with energy $E_{ij}=D_{\mathrm{KL}}(q_i\|\Omega_{ij}q_j)$. Its derivatives apply to this declared energy; replacing it by projected comparison requires the corresponding derivatives of those projected distributions.

We distinguish two objective choices and two derivative-support choices. The canonical objective includes the categorical relative entropy and admits an envelope derivative after optimizing its routing row. The entropy-suppressed surrogate retains the same softmax weights but omits that entropy term from the differentiated scalar. Filtering denotes the receiver's own-row derivative with source distributions held fixed; smoothing includes every row in which the varied state occurs, including later key-side appearances. These distinctions define different vector fields. They do not turn an independently specified predictive content operator into a gradient step.

\subsection{State-Dependent Prior Coupling}
\label{sec:state_dependent_precision}

This auxiliary positive coordinate controls prior anchoring in the belief-refinement objective. It is introduced before the wider variational free energy in which it appears.

\subsubsection{An Auxiliary Positive Weight}

With a fixed prior weight, the self-coupling term in Eq.~\eqref{eq:free_energy_belief} carries a positive constant coefficient. This assumes that prior trust is fixed across contexts.

A more flexible treatment promotes the self-coupling weight from a fixed constant to a per-agent variational parameter $\alpha_i > 0$, subject to a regularizer that prevents degenerate solutions. The derivation below is a single-scale Lagrangian construction and does not invoke any hierarchical or multi-scale structure.

For either the canonical or entropy-suppressed target-blind objective, the adaptive-precision sector is

\begin{equation}
\mathcal{F}_{\alpha,i}
=\alpha_iD_{\mathrm{KL}}(q_i\|p_i)+R(\alpha_i)+\mathcal{A}_i(q,\beta),
\label{eq:free_energy_adaptive}
\end{equation}

where $\mathcal{A}_i$ is the selected attention sector and does not depend explicitly on $\alpha_i$. Both objective paths retain the same precision regularizer. We choose

\begin{equation}
R(\alpha_i) = b_0 \alpha_i - c_0 \log \alpha_i, \qquad b_0 > 0, c_0 > 0,
\label{eq:precision_regularizer}
\end{equation}

For fixed $b_0,c_0$, this regularizer equals the negative log-density of a Gamma distribution with shape $c_0+1$ and rate $b_0$, up to a constant independent of $\alpha_i$ \citep{Bishop2006,murphy2012machine}. The linear term limits large weights and the logarithm excludes zero. This is an auxiliary-weight penalty construction; it does not identify $\alpha_i$ as the precision of a normalized Gaussian observation model.

\subsubsection{Optimal Precision}

Optimizing $\mathcal{F}_{\alpha,i}$ with respect to $\alpha_i$ while holding the beliefs fixed gives

\begin{equation}
\frac{\partial \mathcal{F}_{\alpha,i}}{\partial \alpha_i}
= D_{\mathrm{KL}}(q_i \| p_i) + b_0 - \frac{c_0}{\alpha_i} = 0,
\end{equation}

which gives the optimal state-dependent precision:

\begin{equation}
\boxed{
\alpha_i^*
 = \frac{c_0}{b_0 + D_{\mathrm{KL}}(q_i \| p_i)}.
}
\label{eq:state_dependent_alpha}
\end{equation}

Equivalently, multiplying the Gamma density by the Gibbs factor $\exp(-\alpha_iD)$ gives a normalized Gamma law with rate $b_0+D$ whose mode is the displayed minimizer. This generalized-posterior interpretation does not by itself supply a likelihood for observations with Gaussian precision $\alpha_i$.

When an agent is close to its prior, $\alpha_i^*\approx c_0/b_0$ and the prior resists displacement by alignment. As the full prior KL grows, $\alpha_i^*$ decreases and the alignment sector has greater relative influence.

Writing $D=D_{\mathrm{KL}}(q_i\|p_i)$ gives $\alpha(D)=c_0/(b_0+D)$ and $\alpha'(D)=-c_0/(b_0+D)^2$. On a declared range $0\leq D\leq D_{\max}$, the relative deviation from the zero-divergence value is
\begin{equation}
\frac{\alpha(0)-\alpha(D)}{\alpha(0)}
=\frac{D}{b_0+D}
\leq\frac{D_{\max}}{b_0}.
\end{equation}
Thus $b_0\gg D_{\max}$ gives a quantified near-constant regime, and $c_0/b_0$ is the maximum weight at $D=0$. Learning $b_0,c_0$ under an ordinary probabilistic hierarchy would require its normalization terms. In particular, the Gamma negative log-density also contains $\log\Gamma(c_0+1)-(c_0+1)\log b_0$, which is not constant in those hyperparameters. We do not claim an empirical-Bayes derivation for omitting these terms.

In adaptive mode, $R(\alpha_i)$ remains part of both the canonical and entropy-suppressed objectives. Therefore, for any belief parameter $\theta$, the envelope derivative of the optimized precision sector is

\begin{equation}
\frac{d}{d\theta}
\left[\alpha_i^*D_{\mathrm{KL}}(q_i\|p_i)+R(\alpha_i^*)\right]
=\alpha_i^*\frac{\partial D_{\mathrm{KL}}(q_i\|p_i)}{\partial\theta}.
\label{eq:alpha_envelope_gradient}
\end{equation}

The first derivative contains no explicit $\partial\alpha_i^*/\partial\theta$ term. This cancellation does not extend to the second derivative. Eliminating $\alpha_i$ gives $c_0\log(b_0+D)$ plus a constant; its Hessian includes $-(\alpha_i^*)^2dD\otimes dD/c_0$. Convexity statements below therefore freeze both $\alpha$ and $\beta$, unless their response terms are displayed. For a fixed weight, $\alpha_i$ is constant and the regularizer contributes only an additive constant.


\subsection{Wider Variational Free Energy and the Observation Problem}
\label{sec:final_free_energy}
For a declared observation likelihood, the single-agent free energy is
\begin{equation}
\mathcal F_i^{\mathrm{single}}
=D_{\mathrm{KL}}(q_i\|p_i)-\mathbb E_{q_i}\log p(o_i\mid k_i).
\label{eq:single_agent_fep}
\end{equation}
This motivates combining prior anchoring, communication, and observation terms \citep{Friston2010,Blei2017}.
The belief-refinement family combines this prior coefficient with an attention objective. Section~\ref{sec:supp_full_vfe} gives the complete formulas and their entropy-suppressed comparison. Subsequent sections derive the filtering and smoothing derivatives and natural-gradient updates. These constructions are distinct from the predictive content operator in the main paper.

For next-token prediction, position $i$ uses $x_{\le i}$ and predicts $x_{i+1}$. A target-dependent likelihood can define parameter learning for the composed predictor; placing that held-out target inside its recognition computation would change the information available at prediction time. Evaluating cross-entropy at a mean also differs in general from an expected negative log-likelihood under a Gaussian state. A complete variational account would require a normalized observation model and an inference family that jointly specify these operations. The categorical attention coordinate alone supplies neither that population model nor a descent theorem for a composed predictive architecture.

\paragraph{Interpretation and readout scope.}
\label{sec:symmetry_breaking}
A prior term anchors a state to local reference information. The comparison divergence measures agreement in a declared interface; directions invisible to that interface remain unconstrained by that comparison. A fixed coordinate likelihood generally retains only its stabilizer under frame changes. These observations establish neither a phase transition nor convergence of the predictive stack.
\subsection{Full Variational Free Energy}
\label{sec:supp_full_vfe}

For the optional full-state belief-revision construction, the next token is withheld. Its objective is therefore target-blind: it contains the self-coupling and transported-belief alignment sectors but not the next-token cross-entropy. Write $D_i=D_{\mathrm{KL}}(q_i\|p_i)$, $E_{ij}=D_{\mathrm{KL}}(q_i\|\Omega_{ij}q_j)$, and let $S_i$ denote the active attention support. With the precision regularizer $R$ of Section~\ref{sec:state_dependent_precision}, the canonical inner objective is

\begin{equation}
\begin{aligned}
\mathcal{F}_{\mathrm{belief}}(q,\alpha,\beta)
&=\sum_i\left[
\alpha_iD_i+R(\alpha_i)
+\sum_{j\in S_i}\beta_{ij}E_{ij}
+\tau\sum_{j\in S_i}\beta_{ij}\log\frac{\beta_{ij}}{\pi_{ij}}
\right].
\end{aligned}
\label{eq:free_energy_belief}
\end{equation}

Minimizing the row simplex variables gives $\beta_{ij}^*=\pi_{ij}\exp(-E_{ij}/\tau)/Z_i$, where $Z_i=\sum_{j\in S_i}\pi_{ij}\exp(-E_{ij}/\tau)$. If the precision is adaptive, minimizing it gives $\alpha_i^*$ in Eq.~\eqref{eq:state_dependent_alpha}; for a fixed prior weight, $\alpha_i^*$ below denotes that constant and $R=0$ is a permissible additive-constant convention. The target-blind reduced objective is

\begin{equation}
\boxed{
\mathcal{F}_{\mathrm{belief,red}}[q]
=\sum_i\left[
\alpha_i^*D_i+R(\alpha_i^*)-\tau\log Z_i
\right].
}
\label{eq:free_energy_belief_red}
\end{equation}

The observation-inclusive functional remains useful as a formal generative-model object,

\begin{equation}
\mathcal{F}_{\mathrm{red}}[q]
=\mathcal{F}_{\mathrm{belief,red}}[q]
-\mathbb{E}_q[\log p(o\mid\{k_i\})],
\label{eq:free_energy_final}
\end{equation}

but it is a different functional from the target-blind belief-revision objective. The divergence and alignment sectors are invariant under consistent passive transformations of states and reference distributions; a fixed observation readout need not be. For a channel of comparison dimension $d_a$, a temperature convention may be written $\tau_a=\kappa_a\sqrt{d_a}$, with $\kappa_a>0$. This scalar choice does not identify communication channels with irreducible representations of the frame group.

For any scalar belief parameter $x$, including one that occurs in several rows, the envelope derivative of the target-blind reduced objective is the global double sum

\begin{equation}
\frac{d\mathcal{F}_{\mathrm{belief,red}}}{dx}
=\sum_i\left[
\alpha_i^*\frac{\partial D_i}{\partial x}
+\sum_{j\in S_i}\beta_{ij}^*\frac{\partial E_{ij}}{\partial x}
\right].
\label{eq:envelope_form}
\end{equation}

Terms independent of $x$ are zero. There are no explicit $\partial\beta^*/\partial x$ or $\partial\alpha^*/\partial x$ terms: row-simplex stationarity cancels the former, and stationarity of $\alpha_iD_i+R(\alpha_i)$ cancels the latter. Adding the observation term in Eq.~\eqref{eq:free_energy_final} adds its ordinary derivative to Eq.~\eqref{eq:envelope_form}.

\paragraph{Canonical and entropy-suppressed objectives.} The entropy-suppressed path retains the same regularized self-coupling but replaces each canonical log-partition term by the attention-weighted energy,

\begin{equation}
\mathcal{J}_{\mathrm{belief,surr}}[q]
=\sum_i\left[
\alpha_i^*D_i+R(\alpha_i^*)
+\sum_{j\in S_i}\beta_{ij}^*(q)E_{ij}(q)
\right].
\label{eq:canonical_surrogate_objectives}
\end{equation}

At $\beta^*$, the canonical row identity is

\begin{equation}
-\tau\log Z_i
=\sum_{j\in S_i}\beta_{ij}^*E_{ij}
+\tau\sum_{j\in S_i}\beta_{ij}^*\log\frac{\beta_{ij}^*}{\pi_{ij}}.
\label{eq:canonical_row_decomposition}
\end{equation}

Consequently, for a parameter confined to one row, the surrogate-minus-canonical gradient is

\begin{equation}
\nabla_x\mathcal{J}_{\mathrm{belief,surr}}
-\nabla_x\mathcal{F}_{\mathrm{belief,red}}
=-\tau^{-1}\operatorname{Cov}_{\beta_i^*}
\left(E_{ij},\frac{\partial E_{ij}}{\partial x}\right).
\label{eq:autograd_envelope_gap}
\end{equation}

The analogous global expression sums this covariance over every row in which $x$ occurs. The two vector fields agree only when those covariance terms vanish. Choosing the canonical objective or the entropy-suppressed surrogate is independent of the filtering-versus-smoothing derivative support described next.


\subsection{Causal recognition and optional derivative support}
\label{sec:filtering_free_energy}

Index a length-$N$ sequence as $x_0,\ldots,x_{N-1}$. The autoregressive next-token factorization is

\begin{equation}
p(x_0,\ldots,x_{N-1})
=p(x_0)\prod_{i=0}^{N-2}p(x_{i+1}\mid x_{\leq i}).
\label{eq:next_token_indexing}
\end{equation}

The belief at position $i$ is initialized from $x_i$, may use the prefix $x_{\leq i}$, and predicts the held-out target $x_{i+1}$. Its causal support is $S_i=\{0,\ldots,i\}$, with $\pi_{ij}=1/(i+1)$ on that support and zero otherwise. This indexing is why the target-blind inner objective in Eq.~\eqref{eq:free_energy_belief_red} cannot contain the cross-entropy for $x_{i+1}$. A filtering recognition update uses only information available in this prefix; a smoothing derivative additionally lets later rows alter an earlier belief through its key-side appearances \citep{chung2015recurrent,fraccaro2016sequential,krishnan2017structured}.

In the optional gradient-based branch, filtering applies the own-row update map for either selected objective and excludes variation through the appearances of $q_i$ as a key in later rows. For the canonical path, its mean alignment partial is
\begin{equation}
\frac{\partial E_{ij}}{\partial \mu_i} = \left(\Omega_{ij}\Sigma_j\Omega_{ij}^\top\right)^{-1}\left(\mu_i - \Omega_{ij}\mu_j\right),
\label{eq:queryside_partial}
\end{equation}
which pulls $\mu_i$ toward the gauge-transported neighbor $\Omega_{ij}\mu_j$, and the self-coupling contribution is $\alpha_i^*\Sigma_{p,i}^{-1}(\mu_i-\mu_{\text{prior}})$ as in Eq.~\eqref{eq:mu_gradient_general}.

This own-row step is an update-map iteration, not an exact coordinate minimization, and it does not imply monotonic descent of a global objective. By contrast, smoothing uses the derivative of the selected global objective. Since $q_i$ is the query in row $i$ and a transported key in later rows $m>i$, the canonical smoothing derivative also contains the column term $\sum_{m>i}\beta_{mi}^*\partial E_{mi}/\partial\mu_i$, with
\begin{equation}
\frac{\partial E_{mi}}{\partial \mu_i} = -\Omega_{mi}^\top \left(\Omega_{mi}\Sigma_i\Omega_{mi}^\top\right)^{-1}\left(\mu_m - \Omega_{mi}\mu_i\right).
\label{eq:keyside_partial}
\end{equation}
The surrogate has the same filtering-versus-smoothing support distinction but also carries the through-softmax covariance response in each included row. Derivative support and objective choice are independent mathematical specifications.


\subsection{Optional gradient-based belief dynamics}
\label{sec:supp_belief_dynamics}

The canonical attention value is $\mathcal A_i^{\mathrm{can}}=-\tau\log Z_i$; the surrogate value is
\begin{equation}
\mathcal A_i^{\mathrm{surr}}=\sum_{j\in S_i}\beta_{ij}^*E_{ij}.
\end{equation}
For either selected objective in this optional branch, a filtering step applies the own-row update map

\begin{equation}
q_i^{(t+1)}
=\operatorname{Retr}_{q_i^{(t)}}\left(-\eta_qG_i^{-1}g_i^{\mathrm{filt}}\right),
\qquad
g_i^{\mathrm{filt}}
=\nabla_{q_i}^{\mathrm{own\ row}}
\left[\alpha_i^*D_i+R(\alpha_i^*)+\mathcal{A}_i\right].
\label{eq:belief_dynamics}
\end{equation}
Here $\mathcal{A}_i$ is chosen independently as $\mathcal{A}_i^{\mathrm{can}}$ or $\mathcal{A}_i^{\mathrm{surr}}$. Repeated filtering E-steps are repeated applications of this update map with attention recomputed each time; they are not asserted to minimize the global objective monotonically. In smoothing mode, $g_i$ is instead the derivative of the selected global objective, including every later row in which $q_i$ appears as a key. Neither mode uses the held-out next-token target during the E-step.

The predictive content operator of main-paper Eq.~\eqref{eq:relational_predictive_residual} is a separate update construction. No identification of that operator with the gradient dynamics above, or with a multiscale coarse-graining, follows from the shared routing weights.


\section{Flat-Frame Specialization and Connection Conventions}
\label{app:finite_geometry}

The main construction assigns a Gaussian state to each token and communicates through factorized matrix links. This section isolates its flat-frame specialization on a finite index set. The state coordinates transform under an invertible local frame action. For $U_i\in\mathrm{GL}(K)$, define $\Omega_{ij}=U_iU_j^{-1}$ and the pushforward
\begin{equation}
(\Omega_{ij})_*\mathcal N(\mu_j,\Sigma_j)
=\mathcal N(\Omega_{ij}\mu_j,\Omega_{ij}\Sigma_j\Omega_{ij}^{\top}).
\end{equation}
Under $U_i'=g_iU_i$, transport obeys $\Omega_{ij}'=g_i\Omega_{ij}g_j^{-1}$. The cocycle identity $\Omega_{ij}\Omega_{jk}=\Omega_{ik}$ follows by cancellation. This finite construction uses no continuum derivative, support integral, or coarse-graining assumption.

For comparison with continuum gauge notation, fix the convention $v_i=U_i v_0$ and $D_i=d+A_i$ acting on local column features. A compatible pure-gauge connection is
\begin{equation}
A_i=-(dU_i)U_i^{-1},\qquad
D_i v_i=U_i dv_0.
\label{eq:supp_feature_connection}
\end{equation}
On an overlap with $v_i=\Omega_{ij}v_j$, covariance $D_i v_i=\Omega_{ij}D_jv_j$ requires
\begin{equation}
A_i=\Omega_{ij}A_j\Omega_{ij}^{-1}-(d\Omega_{ij})\Omega_{ij}^{-1}.
\label{eq:supp_connection_overlap}
\end{equation}
The curvature $F_i=dA_i+A_i\wedge A_i$ vanishes for this pure-gauge choice. The body form $B_i=U_i^{-1}dU_i$ is a different trivialization and obeys $A_i=-U_iB_iU_i^{-1}$; it cannot be inserted unchanged into $D_i$ acting on $v_i$. These are standard connection identities \citep{nakahara2003geometry,frankel2011geometry}. Nontrivial connection holonomy requires links beyond this vertex cocycle. General matrix communication links already furnish a different class of covariant edge fields, developed in the main paper, Section~\ref{sec:projection_link_geometry}; their loop products are not the holonomy of the frame-derived connection.

The unrestricted multivariate Gaussian family is closed under these pushforwards. The KL identity used in the paper concerns a common invertible pushforward of two distributions on the same sample space. The matrix links in the main paper may be singular and therefore are not substituted for an invertible pushforward in this full-state KL identity.
\section{Covariance Dynamics of the Optional Belief-Revision Path}
\label{app:covariance_dynamics}

\subsection{Covariance Gradient of the Entropy-Suppressed Surrogate}

We derive the own-row covariance derivative for the optional full-state, target-blind entropy-suppressed surrogate of Section~\ref{app:optional_belief_revision}. This derivative is not the update executed by predictive content. It retains the regularized prior sector $\alpha_i^*D_i+R(\alpha_i^*)$. Differentiating its attention term through $\beta_{ij}^*$ produces $\sum_j(\partial\beta_{ij}^*/\partial\Sigma_i)E_{ij}$. The canonical reduced objective has no such first-order attention response. The envelope cancellation also removes the explicit first-order $\partial\alpha_i^*/\partial\Sigma_i$ term, but the optimized weight contributes to the second derivative. The curvature analysis below makes both response terms explicit.

The own-row surrogate is

\begin{equation}
\mathcal{J}_{\mathrm{surr},i}
=\alpha_i^*D_{\mathrm{KL}}(q_i\|p_i)
+R(\alpha_i^*)
+\sum_{j\in S_i}\beta_{ij}^*D_{\mathrm{KL}}(q_i\|\Omega_{ij}q_j).
\end{equation}

Here $q_i=\mathcal{N}(\mu_i,\Sigma_i)$, $p_i=\mathcal{N}(\mu_{p,i},\Sigma_{p,i})$, and $\sum_{j\in S_i}\beta_{ij}^*=1$.

\subsubsection{Gaussian KL Divergence under $\mathrm{GL}(K)$ Gauge Transport}

The alignment term $D_{\mathrm{KL}}(q_i \| \Omega_{ij} q_j)$ involves KL divergences between Gaussians related by gauge transport $\Omega_{ij} = e^{\phi_i}e^{-\phi_j} \in \mathrm{GL}^+(K)$. The standard Gaussian KL divergence \citep{Bishop2006,murphy2012machine,CoverThomas2006} is

\begin{align}
D_{\mathrm{KL}} \left(
\mathcal{N}(\mu_1, \Sigma_1)
 \|
\mathcal{N}(\mu_2, \Sigma_2)
\right)
&=
\frac{1}{2}
\Big[
\log \frac{|\Sigma_2|}{|\Sigma_1|}
+ \mathrm{tr} \big(\Sigma_2^{-1}\Sigma_1\big)
\\ &\quad
+ (\mu_2 - \mu_1)^\top
  \Sigma_2^{-1}
  (\mu_2 - \mu_1)
- d
\Big].
\label{eq:gaussian_KL}
\end{align}


Differentiating w.r.t.\ $\Sigma_1$ (holding $\mu_1,\mu_2,\Sigma_2$ fixed), using the standard identities $\partial \log|\Sigma|/\partial \Sigma = \Sigma^{-\top}$ and $\partial \mathrm{tr}(A\Sigma)/\partial \Sigma = A^\top$ \citep{petersen2012matrix}, gives

\begin{equation}
\frac{\partial D_{\mathrm{KL}}}{\partial \Sigma_1}
=\frac{1}{2}\left[
-\Sigma_1^{-1}
+ \Sigma_2^{-1}
\right].
\end{equation}

The push-forward $\Omega_{ij} q_j = \mathcal{N}(\Omega_{ij}\mu_j, \Omega_{ij}\Sigma_j\Omega_{ij}^\top)$ used below is the standard linear-transformation rule for a Gaussian random variable under $\Omega \in \mathrm{GL}(K)$ \citep[\S2.3.3]{Bishop2006}. Applying the derivative to the prior and alignment terms gives

\begin{align}
\frac{\partial}{\partial \Sigma_i}
D_{\mathrm{KL}}(q_i \| p_i)
&=\frac{1}{2}\left[
-\Sigma_i^{-1}
+ \Sigma_{p,i}^{-1}
\right],
\\[6pt]
\frac{\partial}{\partial \Sigma_i}
D_{\mathrm{KL}}(q_i \| \Omega_{ij} q_j)
&=
\frac{1}{2}\left[
-\Sigma_i^{-1}
+ (\Omega_{ij}\Sigma_j\Omega_{ij}^\top)^{-1}
\right].
\end{align}

The target-blind inner objective has no observation term. Differentiating the attention-weighted surrogate through the softmax gives

\begin{equation}
\boxed{
\begin{aligned}
\frac{\partial \mathcal{J}_{\mathrm{surr},i}}{\partial \Sigma_i}
&=\frac{1}{2}
\left[
-(1+\alpha_i^*) \Sigma_i^{-1}
+ \alpha_i^*\Sigma_{p,i}^{-1}
+ \sum_{j\in S_i} \beta_{ij}^*
(\Omega_{ij}\Sigma_j\Omega_{ij}^\top)^{-1}
\right] \\
&\quad + \sum_{j\in S_i}
\frac{\partial \beta_{ij}^*}{\partial \Sigma_i}E_{ij}.
\end{aligned}
}
\label{eq:Sigma_gradient_final}
\end{equation}

The first line is the direct derivative shared with the canonical envelope objective. The second line is the surrogate-only response. For two active keys, define $\Delta=E_1-E_2$ and write $E_r'=\partial E_r/\partial\Sigma_i$. The response is exactly

\begin{equation}
-\tau^{-1}\operatorname{Cov}_{\beta}(E,E')
=-\beta_1\beta_2\frac{\Delta}{\tau}(E_1'-E_2').
\label{eq:two_key_near_tie_response}
\end{equation}

For a fixed nonzero gap $|\Delta|=\delta$ and bounded derivatives, its norm is $O(e^{-\delta/\tau}/\tau)$. Along a near-tie path $\Delta=c\tau$, it approaches

\begin{equation}
-\frac{c}{4\cosh^2(c/2)}(E_1'-E_2'),
\label{eq:near_tie_scaled_limit}
\end{equation}

which is generally finite and nonzero rather than $O(\tau)$. At exact equality of the energies the response is zero.

\subsection{Fixed-Point Equation and Symmetric Solution}

For the canonical own-row objective, or for a surrogate configuration in which the response vanishes, setting the direct derivative to zero at $\alpha_i^*=1$ gives

\begin{equation}
\Sigma_i^{-1}
=\frac{1}{2}
\left[
\Sigma_{p,i}^{-1}
+ \sum_{j\in S_i} \beta_{ij}^*
(\Omega_{ij}\Sigma_j\Omega_{ij}^\top)^{-1}
\right].
\label{eq:sigma_fixed_point_beta}
\end{equation}

This is a matrix-valued own-row fixed-point equation. On the entropy-suppressed path, the response in Eq.~\eqref{eq:Sigma_gradient_final} supplies an additional intermediate-temperature term; the canonical first derivative does not contain it.

\subsubsection{Homogeneous limit}

In the homogeneous limit we assume

(i) all agents are identical, so $\Sigma_i = \Sigma_\infty$ for all $i$,

(ii) $\Omega_{ij} \approx I$ (weak misalignment),

(iii) $\Sigma_{p,i} = \Sigma_0$ (shared prior).

Then \eqref{eq:sigma_fixed_point_beta} becomes

\begin{equation}
\Sigma_\infty^{-1}
=\frac{1}{2}
\left[
\Sigma_0^{-1}
+ \Sigma_\infty^{-1}
\right]
\quad\Longrightarrow\quad
\Sigma_\infty = \Sigma_0.
\end{equation}

Hence, in a perfectly symmetric population, the equilibrium covariance reproduces the shared prior.

\subsubsection{Alignment-dominated regime}

Although $\sum_{j\in S_i}\beta_{ij}^*=1$, the concentration of alignment is controlled by $\tau$ in

\begin{equation}
\beta_{ij}^*
=\frac{
\exp \left[
-\frac{1}{\tau}
D_{\mathrm{KL}} \big(q_i \| \Omega_{ij} q_j\big)
\right]
}{\sum_{k\in S_i}
\exp \left[
-\frac{1}{\tau}
D_{\mathrm{KL}} \big(q_i \| \Omega_{ik} q_k\big)
\right]}.
\end{equation}

As $\tau\to0$, $\beta_{ij}^*$ concentrates on the lowest-energy neighbor. Small $\tau$ does not change the coefficient on $\Sigma_i^{-1}$ in Eq.~\eqref{eq:sigma_fixed_point_beta}. A unit coefficient on the alignment term additionally requires $\alpha_i^*\ll1$. In that alignment-dominated regime, the direct fixed point becomes

\begin{equation}
\boxed{
\Sigma_i^{-1}
 \approx 
\sum_{j\in S_i} \beta_{ij}^*
(\Omega_{ij}\Sigma_j\Omega_{ij}^\top)^{-1}
 = \big\langle(\Omega_{ij}\Sigma_j\Omega_{ij}^\top)^{-1}
\big\rangle_{\beta}.}
\label{eq:beta_weighted_precision}
\end{equation}

Here $\langle\cdot\rangle_\beta$ denotes a $\beta^*$-weighted expectation over the active support.

In this regime, agent $i$'s precision matrix becomes the $\beta$-weighted average of transported neighbor precisions. Small $\alpha_i^*$ suppresses the prior term, while small $\tau$ concentrates attention. At $\alpha_i^*=1$, the prior remains at coefficient $1/2$ and Eq.~\eqref{eq:sigma_fixed_point_beta} applies.

When, additionally, all agents already have approximately equal covariances and the transports are near-identity, i.e. $\Sigma_i \approx \Sigma_j$ and $\Omega_{ij} \approx I$, then \eqref{eq:beta_weighted_precision} implies

\begin{equation}
\Sigma_i^{-1} \approx \Sigma_j^{-1}
\quad\Longrightarrow\quad
\Sigma_i \approx \Omega_{ij}\Sigma_j\Omega_{ij}^\top,
\end{equation}

showing that this fixed point is consistent with (rather than deriving) the conditional alignment assumption, since the premises $\Sigma_i \approx \Sigma_j$ and $\Omega_{ij} \approx I$ already presuppose near-alignment.

\subsubsection{Gradient flow dynamics}

Let $\mathcal{O}_i$ denote the selected canonical or surrogate own-row objective. Its covariance flow is

\begin{equation}
\frac{d\Sigma_i}{dt}
=-\eta_\Sigma
\frac{\partial \mathcal{O}_i}{\partial \Sigma_i},
\qquad\eta_\Sigma > 0.
\end{equation}

Fix the means, source covariances, prior covariance, and positive weights $\alpha_i,\beta_{ij}$, with $\sum_j\beta_{ij}=1$. The direct covariance objective is strictly convex on $\mathrm{SPD}(K)$. For a symmetric nonzero direction $H$,
\begin{equation}
d^2\mathcal O_i[H,H]
=\frac{\alpha_i+1}{2}
\left\|\Sigma_i^{-1/2}H\Sigma_i^{-1/2}\right\|_F^2>0.
\label{eq:supp_frozen_covariance_curvature}
\end{equation}
Thus the displayed stationary precision gives the unique minimum of this fixed-weight covariance subproblem. Both weights are frozen in this statement.

Eliminating the prior weight from $\alpha D+b_0\alpha-c_0\log\alpha$ gives $c_0\log(b_0+D)$ plus a constant. Its second variation is
\begin{equation}
d^2\mathcal O_{\mathrm{prior}}[H,H]
=\alpha^*d^2D[H,H]
-\frac{(\alpha^*)^2}{c_0}\bigl(dD[H]\bigr)^2.
\label{eq:supp_alpha_curvature}
\end{equation}
For the canonical attention objective with fixed source prior and temperature, eliminating $\beta$ likewise gives
\begin{equation}
d^2\mathcal O_{\mathrm{align}}[H,H]
=\sum_j\beta_jd^2E_j[H,H]
-\frac{1}{\tau}\operatorname{Var}_{\beta}\bigl(dE_j[H]\bigr).
\label{eq:supp_beta_curvature}
\end{equation}
Consequently the jointly reduced covariance objective need not be convex even if only one of the response terms is present.

A scalar example isolates the prior-weight response. Let the query variance be $s$, the prior variance be $1$, and a single source variance be $3$, with all means zero and $\beta=1$. Define
\begin{equation}
D_1(s)=\tfrac12(s-1-\log s),\qquad
D_3(s)=\tfrac12(s/3-1+\log(3/s)).
\end{equation}
Set $b_0=1/1000$, $c_0=251/1000-\tfrac12\log(3/2)$, and
$J(s)=c_0\log(b_0+D_1(s))+D_3(s)$. At $s_*=3/2$, $\alpha^*=1$, $J'(s_*)=0$, and
\begin{equation}
J''(s_*)=\frac49-\frac{1}{36c_0}\approx-0.13105.
\end{equation}
Freezing $\alpha=1$ instead gives curvature $4/9$. Attention is fixed throughout this counterexample.

Concentrated routing also leaves the prior in the stationary precision. For one source of variance $4$, prior variance $1$, and $\alpha=1$, the stationary query variance is $8/5$. Pairwise alignment requires negligible prior coupling or a prior aligned at the same point. Joint stationarity and local attractivity of the coupled updates require additional conditions and do not follow from fixed-weight convexity.


\section{Gauge Frame Gradients}
\label{app:phi_gradients}

We derive gauge-frame derivatives for the optional full-state canonical target-blind objective and its entropy-suppressed surrogate. These derivatives require the differential of the matrix exponential map and differ for compact and general gauge groups.

\subsection{Structure of the Gauge Frame Gradient}

\paragraph{Convention.} Throughout this section and Section~\ref{app:variational_descent}, the generator-basis components $\partial\mathcal F/\partial\phi_i^a$ use the right-trivialized differential $D_{\phi_i}(\exp)[T_a]=Q_a^{R,(i)}\exp(\phi_i)$, where $Q_a^{R,(i)}=\mathrm{dexp}_{\phi_i}(T_a)$. Thus $\partial\Omega_{ij}/\partial\phi_i^a=Q_a^{R,(i)}\Omega_{ij}$. The Fr\'echet derivative determines the trivialized generator through $Q_a^R=D_\phi(\exp)[T_a]e^{-\phi}$. Equivalently, $D_\phi(\exp)[T_a]=e^\phi Q_a^L$, where $Q_a^L=\mathrm{Ad}_{e^{-\phi}}Q_a^R$. Consistent conversion between these tangent representations leaves the coordinate derivative unchanged. A chosen positive-definite coordinate metric may precondition these components; Section~\ref{sec:glk_preconditioning} distinguishes such choices from a full-$\mathrm{GL}(K)$-invariant metric.

The gauge frames enter through $\Omega_{ij}=\exp(\phi_i)\exp(-\phi_j)$. The global smoothing derivative of the entropy-suppressed surrogate, obtained by differentiating the attention-weighted energies through the softmax, is

\begin{equation}
\frac{\partial \mathcal{J}_{\mathrm{belief,surr}}}{\partial \phi_i}
= \sum_j \left[
\frac{\partial \beta_{ij}}{\partial \phi_i} K_{ij}
+ \beta_{ij} \frac{\partial K_{ij}}{\partial \phi_i}
\right]
+ \sum_k \left[
\beta_{ki} \frac{\partial K_{ki}}{\partial \phi_i}
+ \sum_l \frac{\partial \beta_{kl}}{\partial \phi_i} K_{kl}
\right],
\label{eq:phi_grad_complete}
\end{equation}

where $K_{ij}=D_{\mathrm{KL}}(q_i\|\Omega_{ij}q_j)$. The first sum is the own-row contribution; the second collects key-side contributions from all other rows. The corresponding canonical global smoothing derivative contains only the direct terms,

\begin{equation}
\frac{d\mathcal{F}_{\mathrm{belief,red}}}{d\phi_i}
=\sum_j\beta_{ij}^*\frac{\partial K_{ij}}{\partial\phi_i}
+\sum_k\beta_{ki}^*\frac{\partial K_{ki}}{\partial\phi_i}.
\label{eq:phi_grad_canonical}
\end{equation}

The softmax-response terms in Eq.~\eqref{eq:phi_grad_complete} therefore belong only to the entropy-suppressed surrogate. Independently, filtering retains only the own-row contribution and smoothing retains both sums. Since only $K_{ki}$ depends on $\phi_i$ in row $k$, the surrogate key-side response simplifies to

\begin{equation}
\sum_l \frac{\partial \beta_{kl}}{\partial \phi_i} K_{kl}
= -\frac{\beta_{ki}}{\tau}\frac{\partial K_{ki}}{\partial \phi_i}
\left(K_{ki} - \sum_l \beta_{kl} K_{kl}\right).
\label{eq:reverse_beta_grad_phi}
\end{equation}

This correction vanishes when $K_{ki}$ equals agent $k$'s average divergence $\bar{K}_k = \sum_l \beta_{kl} K_{kl}$, or when $\beta_{ki} \to 0$ (agent $k$ does not attend to $i$).

The coupling weight gradient in the first sum follows from the softmax structure:

\begin{equation}
\frac{\partial \beta_{ij}}{\partial \phi_i} = -\frac{\beta_{ij}}{\tau}\left[
\frac{\partial K_{ij}}{\partial \phi_i} - \sum_k \beta_{ik} \frac{\partial K_{ik}}{\partial \phi_i}
\right].
\label{eq:beta_grad_phi}
\end{equation}

\subsection{Differential of the Matrix Exponential}

All KL gradient terms require the derivative of the transport operator with respect to the gauge frame. For $\Omega_{ij} = \exp(\phi_i)\exp(-\phi_j)$:

\begin{align}
\frac{\partial \Omega_{ij}}{\partial \phi_i}[\xi]
&= \mathrm{dexp}_{\phi_i}(\xi)\exp(\phi_i)\exp(-\phi_j), \label{eq:dOmega_dphi_i}\\
\frac{\partial \Omega_{ij}}{\partial \phi_j}[\xi]
&= -\exp(\phi_i)\mathrm{dexp}_{-\phi_j}(\xi)\exp(-\phi_j). \label{eq:dOmega_dphi_j}
\end{align}

Here $\mathrm{dexp}_{\phi}$ denotes the right-trivialized differential, rather than the matrix-space Fr\'echet derivative itself. For $\xi\in\mathfrak g$ \citep{Gallier2020,Hall2015},

\begin{equation}
\mathrm{dexp}_\phi(\xi) = \sum_{n=0}^{\infty} \frac{1}{(n+1)!} \mathrm{ad}_\phi^n(\xi) = \frac{e^{\mathrm{ad}_\phi} - I}{\mathrm{ad}_\phi}(\xi),
\label{eq:dexp_series}
\end{equation}

where $\mathrm{ad}_\phi(\xi) = [\phi, \xi] = \phi\xi - \xi\phi$ is the adjoint representation. Equation \eqref{eq:dexp_series} defines the right-trivialized differential, related to the Fr\'{e}chet derivative (directional derivative in matrix space) by $D_\phi(\exp)[\xi] = \mathrm{dexp}_\phi(\xi) \cdot e^\phi$. The $\mathrm{GL}(K)$ specialization below uses the Fr\'{e}chet derivative directly via its integral representation~\citep[Ch.~10]{Higham2008}.

\subsubsection{$\mathrm{SO}(N)$ specialization.}

For $\mathrm{SO}(N)$, the gauge frame is expanded in skew-symmetric generators: $\phi = \sum_a \phi^a T_a$ where $T_a \in \mathfrak{so}(N)$, and the gradient is computed component-wise as $\partial \mathcal{F}/\partial \phi^a$. For $\mathrm{SO}(3)$ with $\|\phi\| = \theta$, the differential admits a closed form via the Rodrigues formula:

\begin{equation}
\mathrm{dexp}_\phi(T_a) = T_a + c_1(\theta)[\phi, T_a] + c_2(\theta)[\phi, [\phi, T_a]],
\label{eq:dexp_so3}
\end{equation}

with coefficients:

\begin{equation}
c_1(\theta) = \frac{1 - \cos\theta}{\theta^2}, \qquad c_2(\theta) = \frac{\theta - \sin\theta}{\theta^3}.
\end{equation}

For $\theta < \epsilon$ (numerically $\epsilon \sim 10^{-4}$), Taylor expansions $c_1 \approx \frac{1}{2} - \frac{\theta^2}{24}$ and $c_2 \approx \frac{1}{6} - \frac{\theta^2}{120}$ provide stable evaluation.

\subsubsection{$\mathrm{GL}(K)$ specialization.}

For $\mathrm{GL}(K)$, the gauge frame $\phi \in \mathfrak{gl}(K) \cong \mathbb{R}^{K \times K}$ is a general $K \times K$ matrix. The right-trivialized differential of the exponential does not simplify to a Rodrigues-like formula. The full Fr\'{e}chet derivative $D_\phi(\exp)[\xi] = \mathrm{dexp}_\phi(\xi) \cdot e^\phi$ admits an integral representation \citep[Eq.~10.15]{Higham2008}:

\begin{equation}
D_\phi(\exp)[\xi] = \int_0^1 e^{t\phi} \xi e^{(1-t)\phi} dt,
\label{eq:dexp_integral}
\end{equation}

or equivalently through the $(2K \times 2K)$ block matrix identity \citep[Algorithm~10.27]{Higham2008}:

\begin{equation}
\exp \begin{pmatrix} \phi & \xi \\ 0 & \phi \end{pmatrix}
= \begin{pmatrix} e^\phi & D_\phi(\exp)[\xi] \\ 0 & e^\phi \end{pmatrix}.
\label{eq:dexp_block}
\end{equation}

These give the full Fr\'{e}chet derivative directly; the right-trivialized generator $\mathrm{dexp}_\phi(\xi)$ of Eq.~\eqref{eq:dexp_series} is expressed as $\mathrm{dexp}_\phi(\xi) = D_\phi(\exp)[\xi] \cdot e^{-\phi}$. Either Eq.~\eqref{eq:dexp_integral} or Eq.~\eqref{eq:dexp_block} evaluates the same matrix-space differential.

\subsection{KL Gradient Through Transport}

The KL divergence $K_{ij} = D_{\mathrm{KL}}(q_i \| \Omega_{ij} q_j)$ between $q_i = \mathcal{N}(\mu_i, \Sigma_i)$ and the transported belief $\Omega_{ij} q_j = \mathcal{N}(\Omega_{ij}\mu_j, \Omega_{ij}\Sigma_j\Omega_{ij}^\top)$ depends on $\phi_i$ through $\Omega_{ij}$. Using the Gaussian KL \eqref{eq:gaussian_KL} and the chain rule:

\begin{align}
\frac{\partial K_{ij}}{\partial \phi_i^a}
&= \frac{\partial}{\partial \phi_i^a}\left\{
\frac{1}{2}\left[
\log\frac{|\Omega_{ij}\Sigma_j\Omega_{ij}^\top|}{|\Sigma_i|}
+ \mathrm{tr} \left((\Omega_{ij}\Sigma_j\Omega_{ij}^\top)^{-1}\Sigma_i\right)
+ \Delta\mu_{ij}^\top (\Omega_{ij}\Sigma_j\Omega_{ij}^\top)^{-1} \Delta\mu_{ij}
- K
\right]\right\},
\label{eq:dKL_dphi_expanded}
\end{align}

where $\Delta\mu_{ij} = \mu_i - \Omega_{ij}\mu_j$. Let $\tilde{\Sigma}_{ij} \equiv \Omega_{ij}\Sigma_j\Omega_{ij}^\top$ and $\tilde{\mu}_{ij} \equiv \Omega_{ij}\mu_j$ denote the transported statistics. The derivative has three contributions:

\paragraph{Mean term.} From the Mahalanobis distance:

\begin{equation}
\frac{\partial}{\partial \phi_i^a}\left[(\mu_i - \tilde{\mu}_{ij})^\top \tilde{\Sigma}_{ij}^{-1}(\mu_i - \tilde{\mu}_{ij})\right]
= -2(\mu_i - \tilde{\mu}_{ij})^\top \tilde{\Sigma}_{ij}^{-1} \frac{\partial \tilde{\mu}_{ij}}{\partial \phi_i^a}
- (\mu_i - \tilde{\mu}_{ij})^\top \tilde{\Sigma}_{ij}^{-1} \frac{\partial \tilde{\Sigma}_{ij}}{\partial \phi_i^a} \tilde{\Sigma}_{ij}^{-1} (\mu_i - \tilde{\mu}_{ij}).
\label{eq:mean_term_phi}
\end{equation}

The second term uses the inverse-derivative identity $\partial\tilde{\Sigma}_{ij}^{-1}/\partial\phi_i^a = -\tilde{\Sigma}_{ij}^{-1}(\partial\tilde{\Sigma}_{ij}/\partial\phi_i^a)\tilde{\Sigma}_{ij}^{-1}$, with $\partial\tilde{\Sigma}_{ij}/\partial\phi_i^a$ given by Eq.~\eqref{eq:dtilde_Sigma}; it is generally nonzero and vanishes only where the transported covariance is stationary along $\phi_i^a$.

\paragraph{Trace term.} From the precision-covariance product:

\begin{equation}
\frac{\partial}{\partial \phi_i^a}\mathrm{tr}(\tilde{\Sigma}_{ij}^{-1}\Sigma_i)
= -\mathrm{tr} \left(\tilde{\Sigma}_{ij}^{-1} \frac{\partial \tilde{\Sigma}_{ij}}{\partial \phi_i^a} \tilde{\Sigma}_{ij}^{-1} \Sigma_i\right).
\label{eq:trace_term_phi}
\end{equation}

\paragraph{Log-determinant term.} From the log-determinant ratio:

\begin{equation}
\frac{\partial}{\partial \phi_i^a}\log|\tilde{\Sigma}_{ij}|
= \mathrm{tr} \left(\tilde{\Sigma}_{ij}^{-1} \frac{\partial \tilde{\Sigma}_{ij}}{\partial \phi_i^a}\right).
\label{eq:logdet_term_phi}
\end{equation}

\paragraph{Transport derivatives.} The derivatives of the transported statistics with respect to $\phi_i^a$ are:

\begin{align}
\frac{\partial \tilde{\mu}_{ij}}{\partial \phi_i^a} &= \frac{\partial (\Omega_{ij}\mu_j)}{\partial \phi_i^a}
= Q_a^{(i)} \Omega_{ij} \mu_j, \label{eq:dtilde_mu} \\
\frac{\partial \tilde{\Sigma}_{ij}}{\partial \phi_i^a} &=
Q_a^{(i)} \Omega_{ij} \Sigma_j \Omega_{ij}^\top
+ \Omega_{ij} \Sigma_j \Omega_{ij}^\top (Q_a^{(i)})^\top, \label{eq:dtilde_Sigma}
\end{align}

where $Q_a^{(i)} \equiv \mathrm{dexp}_{\phi_i}(T_a)$ is the right-trivialized differential of the exponential map at $\phi_i$ in the direction of generator $T_a$, defined by the identity $D_{\phi_i}(\exp)[T_a] = Q_a^{(i)} \exp(\phi_i)$. The chain rule applied to $\Omega_{ij} = \exp(\phi_i)\exp(-\phi_j)$ then gives $\partial \Omega_{ij}/\partial \phi_i^a = D_{\phi_i}(\exp)[T_a] \exp(-\phi_j) = Q_a^{(i)} \Omega_{ij}$, which is the form used above. An equivalent expression in terms of the Fr\'echet derivative is $\partial \Omega_{ij}/\partial \phi_i^a = D_{\phi_i}(\exp)[T_a] \exp(-\phi_j)$. The right-trivialized convention is used consistently when the optional frame preconditioner is evaluated; it does not imply that either outer optimizer routes through a retraction dispatcher.

\subsection{Retraction and Numerical Considerations}

\paragraph{Lie algebra updates.} Since $\phi_i \in \mathfrak{g}$ (a vector space), gradient descent proceeds in the Lie algebra directly:

\begin{equation}
\phi_i \leftarrow \phi_i - \eta_\phi \frac{\partial \mathcal{F}}{\partial \phi_i},
\label{eq:phi_update}
\end{equation}

where $\eta_\phi>0$ is the coordinate step size. This displayed Euclidean direction is a coordinate reference, not an invariant natural gradient. The exponential map supplies local group coordinates; choosing a metric on those coordinates or on the group determines a preconditioned direction. Section~\ref{sec:glk_preconditioning} gives explicit choices and their invariance limits.

\paragraph{$\mathrm{SO}(N)$ retraction.} A skew-symmetric tangent $X$ can update a group element by $U^+=\exp(\eta X)U$, which remains orthogonal with determinant one in exact arithmetic. If exponential coordinates are retained, a Baker--Campbell--Hausdorff expansion gives a local approximation to this group product; a finite truncation is not the exact composition. A coordinate norm bound and a group metric bound are different constraints.

\paragraph{$\mathrm{GL}(K)$ updates.} A single real matrix exponential is not surjective onto $\mathrm{GL}^+(K)$ \citep{culver1966existence}. Nevertheless, a product $\exp(\phi_i)\exp(-\phi_j)$ can realize any one prescribed element $A\in\mathrm{GL}^+(K)$: write its polar decomposition as $A=UP$, with $U\in\mathrm{SO}(K)$ and $P\succ0$, choose a skew-symmetric logarithm of $U$ and a symmetric logarithm of $P$, and set $\exp(\phi_i)=U$, $\exp(-\phi_j)=P$ \citep[\S8]{Higham2008}. This is a per-edge existence statement. Shared vertex frames still force $\Omega_{ij}\Omega_{jk}=\Omega_{ik}$, so they do not parameterize independently chosen edge links.

\subsection{Gauge Frame Preconditioning for $\mathrm{GL}(K)$}
\label{sec:glk_preconditioning}

While the basic update~\eqref{eq:phi_update} uses the Euclidean gradient in the Lie algebra, the non-compact structure of $\mathrm{GL}(K)$ introduces a numerical asymmetry that motivates gauge-aware preconditioning. The Lie algebra decomposes via the Cartan involution $\theta(X) = -X^\top$ as:

\begin{equation}
\mathfrak{gl}(K) = \underbrace{\mathfrak{so}(K)}_{\text{compact (skew-symmetric)}} \oplus \underbrace{\mathrm{Sym}_0(K)}_{\text{non-compact (traceless symmetric)}} \oplus \underbrace{\mathbb{R}}_{\text{center (trace)}},
\label{eq:cartan_decomposition}
\end{equation}

where $\mathfrak{so}(K) = \{X : X + X^\top = 0\}$ generates rotations and $\mathrm{Sym}_0(K) = \{X : X = X^\top, \operatorname{tr} X = 0\}$ is the traceless symmetric space generating volume-preserving shears (we reserve $\mathrm{Sym}(K)$ for all symmetric matrices, so that $\mathrm{Sym}(K) = \mathrm{Sym}_0(K) \oplus \mathbb{R}\cdot I$). The right-trivialized factor $\Psi(\operatorname{ad}_\phi)$ distinguishes compact and noncompact adjoint directions: on compact rotations its norm remains bounded away from adjoint resonances, while noncompact adjoint eigenvalue gaps can yield exponential amplification. This statement concerns $\Psi(\operatorname{ad}_\phi)$, not the full Fr\'echet derivative $D_\phi\exp[T]=\Psi(\operatorname{ad}_\phi)(T)\exp(\phi)$. Even when $T$ commutes with $\phi$, the $\Psi$ factor equals the identity but the right factor can grow; for example, $\phi=tI$ gives $D_\phi\exp[T]=e^tT$. The optional preconditioners below address coordinate conditioning without supplying a full-$\mathrm{GL}(K)$ invariant metric.

Four coordinate-conditioning constructions are described below: norm clipping, a Cartan decomposition projector, a regularized Cartan-involution-modified bilinear form, and the Frobenius pullback through the exponential map. Their assumptions and transformation properties differ.

\paragraph{1. Norm clipping (baseline).}
Simple gradient norm clipping $\tilde{g} = g \cdot \min(1, c/\|g\|)$ with threshold $c$. No geometric awareness; treats all Lie algebra directions uniformly.

\paragraph{2. Cartan preconditioning.}
Decompose the gradient into compact and non-compact components and selectively dampen the latter. Given Lie algebra generators $\{T_a\}_{a=1}^{n_{\mathrm{gen}}}$ with Gram matrix $G_{ab} = \operatorname{tr}(T_a^\top T_b)$ and symmetric inner product $S_{ab} = \operatorname{tr}(T_a T_b)$, the projector onto the symmetric subspace is:

\begin{equation}
P_{\mathrm{sym}} = \tfrac{1}{2}G^{-1}(G + S),
\end{equation}

and the preconditioner is $M = I - (1 - \lambda_{\mathrm{sym}}) P_{\mathrm{sym}}$ with dampening factor $\lambda_{\mathrm{sym}} \in (0, 1)$ (typically $0.1$, corresponding to $10\times$ dampening of symmetric directions). The preconditioned gradient is $\tilde{\nabla}_\phi \mathcal{F} = M \cdot \nabla_\phi \mathcal{F}$. Here $P_{\mathrm{sym}}$ acts as $X \mapsto \tfrac{1}{2}(X + X^\top)$ and so retains the trace/dilation direction ($P_{\mathrm{sym}}(I) = I$): it projects onto the full symmetric subspace $\mathrm{Sym}_0(K) \oplus \mathbb{R}\cdot I$ and damps the shear and dilation directions together, rather than the traceless $\mathrm{Sym}_0(K)$ of Eq.~\eqref{eq:cartan_decomposition} alone.

\paragraph{3. Cartan-involution-modified bilinear form (compact-form metric).}
The Killing form of $\mathfrak{gl}(K)$ is $B(X, Y) = 2K \operatorname{tr}(XY) - 2\operatorname{tr}(X)\operatorname{tr}(Y)$. On the trace-free subalgebra $\mathfrak{sl}(K) = \ker(\operatorname{tr})$ it reduces to $B|_{\mathfrak{sl}(K)}(X,Y) = 2K\operatorname{tr}(XY)$, and on the center $\mathbb{R}\cdot I$ it vanishes identically. The Killing form is \emph{indefinite} on real $\mathfrak{sl}(K)$ for $K\geq 2$: under the Cartan decomposition $\mathfrak{sl}(K) = \mathfrak{so}(K) \oplus \mathrm{Sym}_0(K)$ (skew $\oplus$ traceless-symmetric), $B$ is \emph{negative} on $\mathfrak{so}(K)$ and \emph{positive} on $\mathrm{Sym}_0(K)$~\citep{Hall2015,Gallier2020}.

The optional constant construction is not the Killing form itself but the Cartan-involution-modified bilinear form $g(X, Y) = -B(X, \theta(Y))$ with $\theta(X) = -X^\top$ \citep[Ch.~VI]{Knapp2002}. Equivalently, in the generator basis:

\begin{equation}
\tilde{g}_{ab} = 2K \operatorname{tr}(T_a^\top T_b) - 2\operatorname{tr}(T_a)\operatorname{tr}(T_b),
\label{eq:killing_metric}
\end{equation}

which is positive definite on $\mathfrak{sl}(K)$ (the transpose flips the sign on $\mathfrak{so}(K)$, since $\operatorname{tr}(T_a^\top T_a) = -\operatorname{tr}(T_a^2) = \|T_a\|_F^2 > 0$ for skew $T_a$) and degenerate on the center $\mathbb{R} \cdot I$. The regularization $\tilde{g} \mapsto \tilde{g} + \epsilon I$ makes the coordinate matrix positive definite. This compact-form construction agrees with $-B$ up to scale on $\mathrm{SO}(K)$ and assigns positive lengths to both rotations and shears. The Frobenius pairing $\operatorname{tr}(X^\top Y)$ satisfies $\operatorname{tr}((gXg^{-1})^\top(gYg^{-1})) = \operatorname{tr}(X^\top Y)$ for all $X,Y$ when $g^\top g$ is a scalar multiple of the identity. It is therefore invariant under the conformal-orthogonal subgroup $\mathbb{R}^{\times}\mathrm{O}(K)$, whose central scalar factor acts trivially under $\mathrm{Ad}$, but not under a general $\mathrm{GL}(K)$ element. Consequently, $\tilde{g}^{-1}\nabla_\phi\mathcal{F}$ is an $\mathrm{O}(K)$-invariant, more precisely conformal-orthogonal-$\mathrm{Ad}$-invariant, Lie-algebraic preconditioned direction, not a full-$\mathrm{GL}(K)$-covariant natural gradient. A connected Lie group admits a positive-definite bi-invariant metric only when $\mathrm{Ad}(G)$ has compact closure~\citep{Milnor1976}, which fails for noncompact $\mathrm{GL}(K)$.

\paragraph{4. Exponential-coordinate Frobenius pullback (extrinsic and position-dependent).}
This optional construction pulls the Frobenius inner product on $\mathrm{Mat}(K)$ back through the exponential map. The Fr\'echet derivative of $\exp$ at $\phi \in \mathfrak{g}$ in direction $T_a$ admits the right-trivialized form \citep[\S2.7 / Theorem 5.4]{Hall2015}:

\begin{equation}
D_\phi(\exp)[T_a] = \Psi(\mathrm{ad}_\phi)(T_a) \cdot \exp(\phi), \qquad \Psi(z) = \frac{e^z - 1}{z} = \sum_{k=0}^{\infty} \frac{z^k}{(k+1)!},
\end{equation}

where $\mathrm{ad}_\phi(X) = [\phi, X]$ is the adjoint action and $\Psi(z) = (e^z - 1)/z$ is the generating series for the right-trivialized differential; its multiplicative inverse $z/(e^z-1)$ is the Bernoulli generating function whose Taylor coefficients are the Bernoulli numbers and which appears in $d\exp^{-1}$ and the Magnus expansion~\citep{Hall2015}. The equivalent left-trivialized form is $D_\phi(\exp)[T_a] = \exp(\phi) \cdot ((I - e^{-\mathrm{ad}_\phi})/\mathrm{ad}_\phi)(T_a)$, related by the operator identity $\Psi_R(z) = e^z \Psi_L(z)$. The Frobenius pullback metric on $\mathfrak{g}$ is then:

\begin{equation}
\mathcal{G}_{ab}(\phi)
=\langle D_\phi(\exp)[T_a], D_\phi(\exp)[T_b]\rangle_F
= \operatorname{tr}\bigl(\Psi(\mathrm{ad}_\phi)(T_a)^\top \Psi(\mathrm{ad}_\phi)(T_b) \cdot \exp(\phi)\exp(\phi)^\top\bigr),
\label{eq:pullback_metric}
\end{equation}

where $\langle X, Y \rangle_F = \operatorname{tr}(X^\top Y)$ is the Frobenius inner product on $\mathrm{Mat}(K)$. At $\phi = 0$, $\exp(\phi)\exp(\phi)^\top = I$ and $\mathcal{G}_{ab}(0) = G_{ab}$ (the Gram matrix). On the compact subalgebra $\mathfrak{so}(K) \subset \mathfrak{gl}(K)$, $\exp(\phi) \in \mathrm{SO}(K)$ so $\exp(\phi)\exp(\phi)^\top = I$ and the metric reduces to $\langle \Psi(\mathrm{ad}_\phi)(T_a), \Psi(\mathrm{ad}_\phi)(T_b)\rangle_F$. For non-compact (symmetric) $\phi$, $\exp(\phi)\exp(\phi)^\top = \exp(2\phi)$ grows exponentially and conditions the amplification of $D\exp$ in the non-compact sector. This pullback is positive definite exactly where $D_\phi\exp$ has full rank and degenerates at adjoint resonances, where $\Psi(\mathrm{ad}_\phi)$ is singular. It is extrinsic because the ambient Frobenius norm is not invariant under general left multiplication: for $h=\operatorname{diag}(2,1)$, left multiplication sends $E_{11}$ to $hE_{11}$ and changes its squared norm from $1$ to $4$. Thus $\mathcal{G}(\phi)^{-1}\nabla_\phi\mathcal{F}$ is a coordinate preconditioned direction rather than a full-$\mathrm{GL}(K)$-covariant natural gradient. Near rank loss the solve should be regularized as $\mathcal{G} + \epsilon I$. A direct group-element metric such as $g_U(\xi,\eta)=\operatorname{tr}[(U^{-1}\xi)^\top(U^{-1}\eta)]$ is left-invariant, but it would require a group-element parameterization or atlas and is a different geometric construction.

For $n_{\mathrm{gen}}$ generator coordinates, clipping costs $O(n_{\mathrm{gen}})$. Applying a preassembled dense Cartan projector or a precomputed factorization of a constant metric costs $O(n_{\mathrm{gen}}^2)$ per direction. Solving a new dense pullback system costs $O(n_{\mathrm{gen}}^3)$, in addition to assembling the metric. These operation counts do not establish a convergence advantage.


\section{Optional Variational Gradient Descent and Numerical Methods}
\label{app:variational_descent}


\subsection{Natural Gradient Descent on the Gaussian Manifold}
Natural-gradient descent uses the information geometry of the Gaussian family \citep{amari1998natural,Cencov1982,Amari2016}. In mean and covariance coordinates, applying the inverse Fisher metric to coordinate derivatives gives the tangent directions
\begin{equation}
\delta\mu=-\Sigma\nabla_\mu\mathcal F,\qquad
\delta\Sigma=-2\Sigma(\partial_\Sigma\mathcal F)\Sigma.
\label{eq:gaussian_fisher_tangents}
\end{equation}
Here the symmetric covariance derivative is paired with a variation by
\(\operatorname{Tr}((\partial_\Sigma\mathcal F)\delta\Sigma)\). The mean and covariance are separate blocks of the Gaussian Fisher metric.


The covariance update $\delta\Sigma$ is then passed to a manifold retraction (Section~\ref{sec:retraction}) that whitens it into the eigenbasis of $\Sigma$:
\begin{equation}
B = \Lambda^{-1/2} (V^\top \delta\Sigma V) \Lambda^{-1/2},
\end{equation}
where $\Sigma=V\Lambda V^\top$. For step size $\eta_\Sigma$, the full-covariance belief-revision path bounds the step-scaled whitened tangent radially in a Frobenius ball:
\begin{equation}
\eta_\Sigma B\leftarrow
\eta_\Sigma B\min\left(1,\frac{\rho_{\mathrm{tr}}}
{\|\eta_\Sigma B\|_F+\epsilon}\right).
\end{equation}
This is a radial matrix-norm bound. A predictive mean readout alone does not prescribe this covariance dynamics.

\subsection{Manifold Retraction for Covariance Matrices}\label{sec:retraction}
The manifold retraction is the process of mapping the tangent vector back to a valid point on the manifold. This is performed using the affine-invariant exponential map on the manifold of symmetric positive-definite matrices \citep{Pennec2006,Bhatia2007,Absil2008}. Working in the eigenbasis of $\Sigma_k$ with eigenvalues $\lambda_j$, the whitened tangent $B$ is diagonalized $B = U\Lambda_B U^\top$ and the retraction is
\begin{equation}
\Sigma_{k+1} = V \Lambda^{1/2} U \exp(\tau \Lambda_B) U^\top \Lambda^{1/2} V^\top,
\end{equation}
Here $\tau>0$ denotes the covariance step size in this section, not the routing temperature. In exact arithmetic, a symmetric tangent and SPD input give an SPD output. Bounding the whitened tangent limits a proposed step. Additional spectral clipping changes the exponential-map update and does not by itself prove convergence or stability of the coupled dynamics.

\paragraph{Positive definiteness and descent are different conditions.}
A scalar Gaussian example separates them. For variance \(c>0\), let
\begin{equation}
F(c)=\frac12(c-1-\log c),\qquad
\dot c=c-c^2.
\label{eq:scalar_fisher_descent}
\end{equation}
The objective is the KL to a unit-variance Gaussian and the displayed flow follows from Eq.~\eqref{eq:gaussian_fisher_tangents}. At \(c=2\), an Euler step of size \(\eta=2\) gives \(c^+=-2\). The exponential covariance step gives \(c^+=2e^{-2}>0\), but direct substitution yields
\begin{equation}
F(2e^{-2})-F(2)=e^{-2}>0.
\label{eq:spd_not_descent}
\end{equation}
Thus an SPD-preserving retraction does not by itself supply objective descent. An acceptance criterion or a justified step bound is an additional numerical condition; this example embeds in a larger Gaussian by leaving all other coordinates unchanged.

\subsection{Gauge Frame Dynamics on $\mathrm{GL}(K)$}
For the gauge frame fields $\phi_i$, which live in the Lie algebra $\mathfrak{gl}(K)$ (or its subalgebra $\mathfrak{so}(3)$ in the fundamental representation), a different geometric structure must be respected. These fields parameterize the local gauge frames as $U_i = \exp(\sum_a \phi_i^a T_a)$, where $\{T_a\}$ are the Lie algebra generators, and the transport operators take the form $\Omega_{ij} = \exp(\phi_i)\exp(-\phi_j)$. Gradients of the alignment energy $\beta_{ij} D_{\mathrm{KL}}(q_i \| \Omega_{ij}[q_j])$ with respect to $\phi_i$ are obtained via the chain rule through the transport operator: the KL gradient with respect to $\Omega$ is contracted against the differential $\partial\Omega_{ij}/\partial\phi_i^a$.

Per the convention paragraph at the head of Section~\ref{app:phi_gradients}, all gradient components $\partial \mathcal{F}/\partial \phi_i^a$ are defined via the \emph{right-trivialized} differential. The equivalent \emph{left-trivialized} expression is
\begin{equation}
\frac{\partial}{\partial \phi^a}\exp \Bigl(\textstyle\sum_b \phi^b T_b\Bigr) = \exp(X)\cdot Q_a^{L},
\end{equation}
with $X=\sum_b\phi^bT_b$ and $Q_a^L=\mathrm{Ad}_{e^{-X}}(Q_a^R)$ \citep{Hall2015}. The identity $Q_a^Re^X=e^XQ_a^L$ gives the same Fr\'echet derivative and therefore the same chart-coordinate gradient $\partial\mathcal F/\partial\phi^a$. Changing the trivialized tangent representation requires its corresponding group-element factor and, if used, a consistently represented tangent metric; it does not change that assembled coordinate derivative. For the fundamental $\mathfrak{so}(3)$ case a closed-form trivialized generator is available; for general matrices the trivialized generator is related to the Fr\'echet derivative by $Q_a^R=D\exp_X[T_a]e^{-X}$.

The coordinate gradient $\partial\mathcal F/\partial\phi^a$ may be preconditioned by a positive-definite coordinate metric, including the regularized Cartan construction or the exponential-coordinate Frobenius pullback of Eq.~\eqref{eq:pullback_metric}. If a group-element update is used instead, its tangent representation and metric must be specified on the group. An additive coordinate step and a retracted group tangent step are different finite-step maps.


\paragraph{Derivative checks.} Central finite differences of a selected scalar objective provide a numerical check of its coordinate derivatives. Comparing only two calculations of the matrix-exponential differential does not check the assembly of the complete frame derivative: the direct and softmax-response terms must be included according to the selected objective. Numerical agreement is a diagnostic and does not replace the analytic derivations above.


\section{Conditional Uniqueness of the Forward KL Divergence via Variational Duality}
\label{sec:supp_conditional_uniqueness}

\paragraph{Intuition.}
Why does the inter-agent term take the forward-KL form rather than another statistical distance? We impose a local convex $f$-divergence that enters linearly in $\beta_{ij}$ and require a log-linear Gibbs stationary update. The envelope theorem then identifies the marginal attention cost with the chosen divergence, but that identity holds for every member of the linearly coupled class and does not select forward KL.

Generic alternatives break this log-linearity. For example, the $\chi^2$ generator $f(t)=(t-1)^2/2$ has $f'(t)=t-1$, so its stationarity equation mixes $q_i$ with $\log q_i$. Log-linear closure instead selects the positive ray $f_c(t)=c(t\log t-t+1)$. The exact half-exponent update used below, or the normalization $f''(1)=1$, fixes the remaining scale to $c=1$.

Within the convex analytic $f$-divergence class, the theorem proves unit-scale uniqueness only under its exact exponent target and a fixed admissible witness whose density ratio has a nonempty open essential range, or under overlap conditions that force witness-specific constants to agree.

In what follows we restrict to the matched-fiber framework, where all agents share the same fiber dimension $K$ and representation $\rho$, so transport acts within one vector space. Extending the theorem to unequal fibers would require specified bundle morphisms that are invertible on the compared measures, preserve the admissible density family, and satisfy the compatibility conditions used by the variational argument; no such extension is proved here. Throughout this section $\Omega_{ij}q_j$ denotes the push-forward density $[(\Omega_{ij})_*q_j]$, including its change-of-variables Jacobian, rather than pointwise multiplication of a matrix by a scalar density.

The witness condition is substantive: the chosen admissible density family must contain the witness and the transformed reference laws used by the argument. Closure under an arbitrary invertible change of variables cannot be assumed for a restricted covariance family.

\subsubsection{The Coupled Variational Problem}

For fixed attention coordinates $\beta_i$, define the belief-coordinate value function

\begin{equation}
F_i(\beta_i)
=\min_{q_i}
\Big\{
D_{\mathrm{KL}}(q_i \| p_i)
+ \sum_{j\neq i} \beta_{ij} C_{ij}(q_i)
\Big\},
\label{eq:F_i_supp}
\end{equation}

where $C_{ij}(q_i):=\mathcal{D}(q_i,\Omega_{ij}q_j)$ and a selected minimizer is denoted $q_i^*(\beta_i)$. This value function is distinct from the attention-coordinate objective evaluated at a fixed belief.

With $H_i(c)=-\log p_i(c)$ and normalized densities, the KL divergence has the exact decomposition

\begin{equation}
D_{\mathrm{KL}}(q_i \| p_i) = \langle H_i \rangle_{q_i} - S(q_i),
\end{equation}

where

\begin{align}
H_i(c) &:= -\log p_i(c) \quad \text{(local "energy")}, \\
S(q_i) &:= -\int q_i(c) \log q_i(c) dc \quad \text{(Shannon entropy)}, \\
\langle H_i \rangle_{q_i} &:= \int q_i(c) H_i(c) dc \quad \text{(expected energy)}.
\end{align}

This yields $D_{\mathrm{KL}}(q_i \| p_i) = \int q_i(c)[\log q_i(c) - \log p_i(c)] dc$.

At a fixed $q_i$, the attention coordinates $\beta_{ij}\ge 0$, $\sum_j\beta_{ij}=1$, minimize

\begin{equation}
\mathcal{J}_i(\beta_i\mid q_i)
=\sum_{j\neq i}\beta_{ij} C_{ij}(q_i)
+ \tau \sum_{j\neq i}\beta_{ij}\log\frac{\beta_{ij}}{\pi_{ij}},
\label{eq:J_i_supp}
\end{equation}

where $\pi_{ij}>0$ is the row prior on the active support. The simplex KKT equation gives
\begin{equation}
C_{ij}(q_i)+\tau\left(\log\frac{\beta_{ij}}{\pi_{ij}}+1\right)+\lambda_i=0,
\end{equation}
and hence the explicit fixed-belief update
\begin{equation}
\beta_{ij}^*(q_i)
=\frac{\pi_{ij}\exp[-C_{ij}(q_i)/\tau]}
{\sum_k\pi_{ik}\exp[-C_{ik}(q_i)/\tau]}.
\label{eq:beta_fixed_q_supp}
\end{equation}
The entropy term and the row prior are part of this coordinate objective; replacing them by $\tau\sum_j\beta_{ij}\log\beta_{ij}$ silently assumes a uniform prior up to a row constant.

\subsubsection{The Forward KL and the Geometric-Mean Solution}
Let $\mathcal{D}(q_i,q_j) = D_{\mathrm{KL}}(q_i \| \Omega_{ij} q_j)$, where

\[
\frac{\delta D_{\mathrm{KL}}(q_i \| \Omega_{ij} q_j)}{\delta q_i(c)}
= \log\frac{q_i(c)}{\Omega_{ij}(c) q_j(c)} + 1.
\]

The stationarity condition for \eqref{eq:F_i_supp} is

\begin{equation}
H_i(c)
+ \log q_i(c)
+ \sum_j \beta_{ij}
 \left[
\log\frac{q_i(c)}{\Omega_{ij}(c)q_j(c)}+1
\right]
= \lambda_i,
\end{equation}

with $\lambda_i$ enforcing normalization. The additive constant $\sum_j \beta_{ij} \cdot 1 = 1$ (using $\sum_j \beta_{ij} = 1$) is independent of $q_i(c)$ and is absorbed into $\lambda_i$; it therefore does not affect $q_i^*$, and the equivalent $f$-divergence form below, which uses $f'(r) = \log r$ with no additive constant, gives the identical stationary point.

Solving for $q_i(c)$ then yields a Boltzmann distribution whose mean field is the geometric average of the transported neighbor beliefs.

\begin{equation}
q_i^*(c)
= \frac{1}{Z_i}
e^{-H_i(c)/2}
\prod_j
[\Omega_{ij}(c) q_j(c)]^{\beta_{ij}/2},
\label{eq:q_i_star_supp}
\end{equation}

This structure is preserved only for the forward KL divergence; within the convex $f$-divergence class assumed here, alternative $f$-divergences destroy the log-linearity of the exponent as previously discussed (integral probability metrics such as Wasserstein distance or MMD lie outside this class and are out of scope by hypothesis rather than ruled out).

\subsubsection{Dual Relation via the Envelope Theorem}
At the belief-coordinate minimizer $q_i^*(\beta_i)$, the envelope theorem implies

\begin{equation}
\frac{\partial F_i}{\partial \beta_{ij}}
= C_{ij}\bigl(q_i^*(\beta_i)\bigr)
= D_{\mathrm{KL}}\bigl(q_i^*(\beta_i) \| \Omega_{ij}q_j\bigr),
\label{eq:envelope_supp}
\end{equation}

for an unconstrained coordinate derivative. On the simplex, the feasible transfer derivative that moves mass from source $k$ to source $j$ is $C_{ij}(q_i^*)-C_{ik}(q_i^*)$; adding one row-constant value to every marginal cost changes neither this derivative nor the softmax.

If belief minimization is substituted first, the fully reduced attention functional is
\begin{equation}
\widetilde{\mathcal J}_i(\beta_i)
=F_i(\beta_i)+\tau\sum_j\beta_{ij}\log\frac{\beta_{ij}}{\pi_{ij}}.
\label{eq:J_reduced_beta_supp}
\end{equation}
Its stationarity equation is the implicit Gibbs fixed point
\begin{equation}
\beta_{ij}
\propto \pi_{ij}\exp\left[-C_{ij}\bigl(q_i^*(\beta_i)\bigr)/\tau\right].
\label{eq:beta_implicit_supp}
\end{equation}
It is not obtained by minimizing $\sum_j\beta_{ij}C_{ij}(q_i^*(\beta_i))$ and differentiating that product, which would double-count the response of the optimized belief. At a joint coordinate fixed point, the fixed-$q_i$ update~\eqref{eq:beta_fixed_q_supp} and the reduced fixed point~\eqref{eq:beta_implicit_supp} coincide.

\subsubsection{Conditional Uniqueness Theorem}
\label{thm:uniqueness_supp}

\textbf{Theorem.} \textit{Let\/ $\mathcal{D}(q,p)$ be a convex $f$-divergence~\citep{Csiszar1967} of the form
\[
\mathcal{D}(q,p)
= \int p(c) f \left(\frac{q(c)}{p(c)}\right)dc,
\]
with $f:(0,\infty)\to\mathbb{R}$ convex, $f(1)=0$, real-analytic on $(0,\infty)$, nondegenerate in the sense that $f''(1)>0$, and, without loss of generality\footnote{Any $f$-divergence is invariant under $f(t)\mapsto f(t)+a(t-1)$ when $q$ and $p$ are normalized, so we may choose the representative with $f'(1)=0$.}, $f'(1)=0$. Suppose $\mathcal{D}$ is linear inside the free energy \eqref{eq:F_i_supp} and $\sum_j\beta_{ij}=1$. Assume that there is one fixed admissible witness configuration $(p_i,\{q_j\},\beta_i)$ and a reference density $\bar q$ such that all transported neighbors $\Omega_{ij}q_j$ equal $\bar q$ almost everywhere and the ratio $r(c)=q_i^*(c)/\bar q(c)$ has an essential range containing a nonempty open interval. Equivalently, a family of witness configurations may be used only when overlaps of their open essential ratio ranges force their configuration-specific residual constants to agree. Then the stationary solution $q_i^*$ assumes the geometric-mean Boltzmann form}

\begin{equation}
q_i^*(c) \propto p_i(c)^{1/2} \prod_{j}\bigl[\Omega_{ij}(c) q_j(c)\bigr]^{\beta_{ij}/2}
\label{eq:geometric_mean_target_supp}
\end{equation}

\textit{for all choices of prior $p_i$ and neighbor beliefs $\{q_j\}$ if and only if $\mathcal{D}$ is the forward Kullback-Leibler divergence, up to the affine equivalence above:}

\[
\mathcal{D}(q,p) = D_{\mathrm{KL}}(q \| p).
\]

\paragraph{Proof.}

\textbf{Step 1. The variational derivative of a general $f$-divergence.}

With $\mathcal{D}(q_i,\Omega_{ij}q_j)=\int (\Omega_{ij}q_j) f \bigl(q_i/(\Omega_{ij}q_j)\bigr) dc$ and ratio $r(c)=q_i(c)/\bigl[\Omega_{ij}(c) q_j(c)\bigr]$, the functional derivative is

\begin{equation}
\frac{\delta\mathcal{D}}{\delta q_i(c)} = f'(r(c)).
\end{equation}

The stationarity condition $\delta F_i/\delta q_i(c)=\lambda$ (with Lagrange multiplier $\lambda$) is then

\begin{equation}
\log\frac{q_i(c)}{p_i(c)} + 1
+ \sum_j \beta_{ij} f' \left(\frac{q_i(c)}{\Omega_{ij}(c) q_j(c)}\right) = \lambda.
\label{eq:general_stationarity_supp}
\end{equation}

\textbf{Step 2. Forward implication ($D_{\mathrm{KL}}\Rightarrow$ geometric mean).}

Next, for $\mathcal{D}=D_{\mathrm{KL}}$, the generator is $f(t)=t\log t - t + 1$, yielding $f'(t)=\log t$. Substituting into \eqref{eq:general_stationarity_supp} we have:

\[
\log q_i - \log p_i + 1 + \sum_j \beta_{ij}\bigl(\log q_i - \log\Omega_{ij}q_j\bigr) = \lambda.
\]

However, since $\sum_j\beta_{ij}=1$, the $\log q_i$ terms become $2\log q_i$:

\[
2\log q_i = \log p_i + \sum_j\beta_{ij}\log(\Omega_{ij}q_j) + (\lambda - 1).
\]

By exponentiating we find the geometric-mean form \eqref{eq:geometric_mean_target_supp}.

\textbf{Step 3. Reverse implication (geometric mean $\Rightarrow D_{\mathrm{KL}}$).}

Assume the minimizer has the geometric mean form \eqref{eq:geometric_mean_target_supp} for all priors $p_i$ and neighbor beliefs $\{q_j\}$.

Taking logarithms yields:

\begin{equation}
2\log q_i^*(c) = \log p_i(c) + \sum_j\beta_{ij}\log\bigl(\Omega_{ij}(c) q_j(c)\bigr) + C.
\label{eq:log_target_supp}
\end{equation}

Next, rearranging and using $\sum_j\beta_{ij}=1$ we have:

\begin{equation}
\log\frac{q_i^*}{p_i} + \sum_j\beta_{ij}\log\frac{q_i^*}{\Omega_{ij}q_j} = C'.
\label{eq:rearranged_target_supp}
\end{equation}

Subtracting \eqref{eq:rearranged_target_supp} from the general stationarity condition \eqref{eq:general_stationarity_supp} gives
\begin{equation}
\sum_{j} \beta_{ij}\Bigl[ f'\left(\frac{q_i^*(c)}{\Omega_{ij}(c)q_j(c)}\right) - \log\frac{q_i^*(c)}{\Omega_{ij}(c)q_j(c)}\Bigr] = k,
\label{eq:residual_sum_supp}
\end{equation}
Here $k=\lambda-1-C'$ is constant in $c$ but can depend on the configuration. Evaluate this equation in the single fixed witness from the theorem. All transported neighbors then equal $\bar q$, all ratios coincide with $r(c)=q_i^*(c)/\bar q(c)$, and $\sum_j\beta_{ij}=1$ reduces the equation to
\begin{equation}
f'(r(c)) - \log r(c) = k,
\label{eq:ode_supp}
\end{equation}
The essential range of $r$ contains an open interval while $k$ remains fixed. Continuity therefore gives $f'(t)-\log t=k$ on that interval, and real analyticity extends the identity to all $t>0$. If several witnesses are used instead, the same conclusion requires overlaps that identify their constants before analytic continuation.

Next, by integrating we have:

\[
f(t) = t\log t - t + kt + C_0.
\]

If we impose $f'(1)=0$ then $k=0$. Similarly by imposing $f(1)=0$ then $C_0 = 1$. Therefore, we have

\[
f(t) = t\log t - t + 1,
\]

which is the generator of $D_{\mathrm{KL}}(q \| p)$. \hfill $\square$

\paragraph{Why pooled configurations are insufficient.}
On a two-point space with reference density $(1/2,1/2)$, normalization gives $r_1+r_2=2$. Consider
\begin{equation}
f(t)=t\log t-t+1+\frac{(t-1)^3}{3}.
\end{equation}
This generator is real analytic and strictly convex on $(0,\infty)$ because
\begin{equation}
f''(t)=\frac{2t^2-2t+1}{t}>0.
\end{equation}
Nevertheless, $f'(t)-\log t=(t-1)^2$, so the two realized ratios in each fixed configuration satisfy $(r_1-1)^2=(r_2-1)^2=k_p$. As the prior varies, the pooled ratios cover $(0,2)$ but $k_p$ varies with the configuration. An open range obtained only by pooling configurations is therefore insufficient. The required fixed witness is also not automatic for diagonal Gaussians under a general transport, because inverse transport generally produces a full covariance outside the diagonal family.

\paragraph{Remark (closure and scale).}
If one requires only log-linear Gibbs closure rather than the exact half-exponent target, the admissible generators form the positive KL ray
\begin{equation}
f_c(t)=c(t\log t-t+1), \qquad c>0.
\end{equation}
With unit weight on the prior term, stationarity gives
\begin{equation}
q_i^*(x)\propto p_i(x)^{1/(1+c)}
\prod_j\bigl[\Omega_{ij}(x)q_j(x)\bigr]^{c\beta_{ij}/(1+c)}.
\end{equation}
Thus closure alone selects $cD_{\mathrm{KL}}$, not its unit scale. The convention $f''(1)=1$ fixes $c=1$, as does the exact target with prior exponent $1/2$ and neighbor exponents $\beta_{ij}/2$.

\paragraph{Remark (general self-coupling weight).} The objective \eqref{eq:F_i_supp} carries unit weight on the prior term $D_{\mathrm{KL}}(q_i \| p_i)$, corresponding to the fixed choice $\alpha_i=1$. For a general self-coupling weight $\alpha_i > 0$, replacing $D_{\mathrm{KL}}(q_i \| p_i)$ by $\alpha_i D_{\mathrm{KL}}(q_i \| p_i)$ (Section~\ref{sec:state_dependent_precision}), the stationarity condition acquires the coefficient $\alpha_i + 1$ on $\log q_i^*$ in place of $2$, so the geometric-mean solution generalizes to
\begin{equation}
q_i^*(c) \propto p_i(c)^{\alpha_i/(\alpha_i + 1)} \prod_{j}\bigl[\Omega_{ij}(c) q_j(c)\bigr]^{\beta_{ij}/(\alpha_i + 1)},
\end{equation}
which reduces to \eqref{eq:geometric_mean_target_supp} at $\alpha_i = 1$. This formula uses the unit-scale generator $c=1$. More generally, $f_c$ gives prior exponent $\alpha_i/(\alpha_i+c)$ and neighbor exponents $c\beta_{ij}/(\alpha_i+c)$; a scale convention or an exact exponent target is therefore required for uniqueness at every $\alpha_i>0$.

\subsubsection{Dual Cost Identification}
Under the theorem's fixed-witness and unit-scale hypotheses, the envelope theorem applied to the belief-coordinate value function gives the marginal attention cost

\begin{equation}
C_{ij}\bigl(q_i^*(\beta_i)\bigr)
\equiv \frac{\partial F_i}{\partial\beta_{ij}}
= D_{\mathrm{KL}}\bigl(q_i^*(\beta_i) \| \Omega_{ij}q_j\bigr),
\end{equation}

up to a common row constant when derivatives are restricted to the simplex. The fixed-belief attention coordinate then uses $C_{ij}(q_i)$ in Eq.~\eqref{eq:beta_fixed_q_supp}; after exact belief minimization the same relation appears as the implicit fixed point~\eqref{eq:beta_implicit_supp}.

\subsubsection{Interpretations}
Gauge covariance fixes the comparison between $q_i$ and the transported $\Omega_{ij}q_j$ in a common frame. The envelope identity $\partial F_i/\partial\beta_{ij}=\mathcal D(q_i^*(\beta_i),\Omega_{ij}q_j)$ holds throughout the linearly coupled class and is therefore a consistency property, not a selection criterion. It must not be confused with the fixed-belief coordinate objective that produces the explicit softmax. Log-linear Gibbs closure selects the positive KL ray; the fixed-witness hypothesis together with either the exact half-exponent target or $f''(1)=1$ selects its unit-scale forward-KL member. The forward KL is also the Bregman divergence generated by negative entropy, whose Hessian induces the Fisher-Rao metric. Its information-geometric role is consistent with the classical results of \citet{Cencov1982} and the standard characterization of forward KL as the intersection of the Bregman and $f$-divergence classes \citep{Amari2016}; \citet{Petz1996} describes the distinct nonunique family of monotone metrics in the quantum setting.

\section{Belief Gradients and Residual Update Maps}\label{app:belief_revision}

This section derives the belief-coordinate vector fields for the full-state objective and distinguishes them from a prescribed predictive content operator. The mean-only subproblem keeps covariance positive and fixed; it does not take a singular zero-variance limit of the complete Gaussian KL. Filtering uses the own-row partial derivative. Smoothing differentiates the global scalar through every query-side and key-side occurrence.

A belief refinement and a parameter-learning step need not optimize the same scalar functional. Interpreting them as variational EM requires a specified normalized model, a variational family, and compatible coordinate objectives; the existence of differentiable update maps alone supplies no EM monotonicity statement \citep{Dempster1977,Neal1998}.

\subsection{Observation Likelihood as Loss Function.}

In the proposed gauge reading, observations act as a source term whose likelihood generally reduces the vacuum reparameterization redundancy to the likelihood stabilizer. The observation-free free energy, written with $\beta_{ij}$ at its optimal softmax value, takes the reduced form (cf.\ Eq.~\ref{eq:free_energy_final}):

\begin{equation}
\mathcal{F}_{\text{vacuum}}[\{q_i\}]
=\sum_i\left[\alpha_i^*D_{\mathrm{KL}}(q_i\|p_i)+R(\alpha_i^*)-\tau\log Z_i\right].
\end{equation}

Introducing per-agent observations $o_i$ adds the likelihood term:
\begin{equation}
\mathcal{F}[\{q_i\}] = \mathcal{F}_{\text{vacuum}}[\{q_i\}] - \sum_i \mathbb{E}_{q_i}[\log p(o_i \mid k_i)].
\end{equation}

If $q_i$ concentrates at $\mu_i$ and the log likelihoods are uniformly integrable under this limit, the observation term satisfies:

\begin{equation}
-\mathbb{E}_{q_i}[\log p(o_i \mid k_i)] \to -\log p(o_i \mid \mu_i),
\end{equation}
the standard negative log-likelihood.

For Gaussian observations $p(o \mid \mu) = \mathcal{N}(o \mid \mu, \Sigma)$:

\begin{equation}
-\log p(o \mid \mu) = \frac{1}{2}(o - \mu)^\top \Sigma^{-1} (o - \mu) + \text{const.}
\end{equation}

\begin{equation}
\mathcal{L}_{\text{obs}}
=\frac{1}{2}\|o-\mu\|^2_{\Sigma^{-1}}
\quad\text{(Mahalanobis squared error)}.
\end{equation}

For categorical observations $p(o \mid \mu) = \mathrm{Categorical}(\mathrm{softmax}(\mu))$:

\begin{equation}
-\log p(o \mid \mu) = -\sum_k o_k \log(\mathrm{softmax}(\mu)_k) \quad \text{(Cross-entropy loss)}.
\end{equation}

These are standard loss functions in machine learning.


\paragraph{Fixed-precision mean-only comparison.}

The singular family $q_i=\mathcal{N}(\mu_i,\epsilon I)$ with $\epsilon\to0^+$ does not define an adaptive-precision limit of the full Gaussian KL: its covariance contribution diverges, and substituting only the finite mean term into Eq.~\eqref{eq:state_dependent_alpha} changes the objective. We therefore do not take that limit. Instead, hold $\alpha_i=\alpha>0$ fixed, hold $q_i=\mathcal{N}(\mu_i,\sigma^2I)$ at a common positive covariance during the mean update, set $\Omega_{ij}=I$, and compare only the mean-dependent terms. Then

\begin{equation}
\alpha D_{\mathrm{KL}}
\left(
\mathcal{N}(\mu_i,\sigma^2I)
\middle\|
\mathcal{N}(\mu_{\mathrm{prior}},\sigma_p^2I)
\right)
=\frac{\alpha}{2\sigma_p^2}
\|\mu_i-\mu_{\mathrm{prior}}\|^2+C_{\mathrm{cov}},
\label{eq:alpha_i_deterministic}
\end{equation}

where $C_{\mathrm{cov}}$ is independent of $\mu_i$. Dropping such constants and defining $E_{ij}^{\mathrm{mean}}=\|\mu_i-\mu_j\|^2/(2\sigma^2)$, the canonical target-blind mean objective is

\begin{equation}
\mathcal{F}_{\mathrm{belief,red}}^{\mathrm{mean}}
=\sum_i\left[
\frac{\alpha}{2\sigma_p^2}\|\mu_i-\mu_{\mathrm{prior}}\|^2
+R(\alpha)-\tau\log Z_i^{\mathrm{mean}}
\right],
\label{eq:deterministic_free_energy_adaptive}
\end{equation}

where $Z_i^{\mathrm{mean}}=\sum_{j\in S_i}\pi_{ij}\exp(-E_{ij}^{\mathrm{mean}}/\tau)$. The entropy-suppressed mean surrogate is

\begin{equation}
\mathcal{J}_{\mathrm{belief,surr}}^{\mathrm{mean}}
=\sum_i\left[
\frac{\alpha}{2\sigma_p^2}\|\mu_i-\mu_{\mathrm{prior}}\|^2
+R(\alpha)
+\sum_{j\in S_i}\beta_{ij}^*E_{ij}^{\mathrm{mean}}
\right].
\label{eq:deterministic_free_energy_standard}
\end{equation}

Because $\alpha$ is fixed, $R(\alpha)$ is constant in this mean-gradient comparison. Neither equation contains the held-out prediction target. A next-token cross-entropy objective may be applied to parameters after the mean update; it is not a term in this target-blind mean subproblem.

\subsection{Gradient Descent Dynamics.}

Let $\delta_i=\mu_i-\mu_{\mathrm{prior}}$, $\delta_{ij}=\mu_i-\mu_j$, and $\lambda_p=\alpha/\sigma_p^2$. The canonical own-row Euclidean derivative of Eq.~\eqref{eq:deterministic_free_energy_adaptive} is

\begin{equation}
g_{i,\mathrm{can}}^{\mathrm{filt}}
=\lambda_p\delta_i
+\frac{1}{\sigma^2}\sum_{j\in S_i}\beta_{ij}^*\delta_{ij}.
\label{eq:isotropic_canonical_gradient}
\end{equation}

For the entropy-suppressed objective, differentiating through the softmax adds

\begin{equation}
g_{i,\mathrm{surr}}^{\mathrm{filt}}
=g_{i,\mathrm{can}}^{\mathrm{filt}}
+\frac{1}{2\sigma^2}
\sum_{j\in S_i}
\frac{\partial\beta_{ij}^*}{\partial\mu_i}
\|\delta_{ij}\|^2.
\label{eq:isotropic_surrogate_gradient}
\end{equation}

The mean-block Fisher information of $\mathcal{N}(\mu_i,\Sigma_i)$ is $G_{\mu_i}=\Sigma_i^{-1}$, so the natural-gradient step uses $G_{\mu_i}^{-1}=\Sigma_i$. For $\Sigma_i=\sigma^2I$, the canonical update with the ordinary step size $\eta$ is

\begin{equation}
\boxed{
\mu_i^{(t+1)}
=\mu_i^{(t)}
-\eta\left[
\frac{\alpha\sigma^2}{\sigma_p^2}\delta_i
+\sum_{j\in S_i}\beta_{ij}^*\delta_{ij}
\right].
}
\label{eq:isotropic_natural_update}
\end{equation}

For the surrogate, the bracket additionally contains $\frac{1}{2}\sum_j(\partial\beta_{ij}^*/\partial\mu_i)\|\delta_{ij}\|^2$. No second factor of $\sigma^2$ is absorbed into $\eta$. The alignment term equals $\mu_i-\sum_j\beta_{ij}^*\mu_j$ because the row weights sum to one.

\paragraph{General (non-isotropic) mean gradient.}
Let $\chi_{\mathrm{surr}}=0$ for the canonical objective and $\chi_{\mathrm{surr}}=1$ for the entropy-suppressed surrogate. The global smoothing derivative with respect to $\mu_i$ is

\begin{equation}
\boxed{
\begin{aligned}
g_i^{\mathrm{smooth}}
&=\alpha_i^*\Sigma_{p,i}^{-1}
(\mu_i-\mu_{\mathrm{prior}})\\
&\quad+\sum_r\sum_{j\in S_r}
\left[
\beta_{rj}^*\frac{\partial E_{rj}}{\partial\mu_i}
+\chi_{\mathrm{surr}}
\frac{\partial\beta_{rj}^*}{\partial\mu_i}E_{rj}
\right].
\end{aligned}
}
\label{eq:mu_gradient_general}
\end{equation}

The double sum includes the own-row query terms and every key-side term from later rows. Filtering retains only $r=i$ and the query-side dependence. The adaptive precision contributes the factor $\alpha_i^*$ but no response term because $R(\alpha_i^*)$ is retained and Eq.~\eqref{eq:alpha_envelope_gradient} applies. The through-softmax term appears only for the entropy-suppressed surrogate. The corresponding natural-gradient step is

\begin{equation}
\boxed{
\mu_i^{(t+1)}
=\operatorname{Retr}_{\mu_i^{(t)}}
\left(-\eta\Sigma_i g_i\right),
}
\label{eq:natural_gradient_general}
\end{equation}

The natural-gradient label here is specific to the belief channel: the $(\mu_i, \Sigma_i)$ updates use the Fisher-Rao metric of the Gaussian manifold, exact for these coordinates. The gauge-frame sector has the optional regularized Cartan/Killing Lie-algebraic conditioner
\begin{equation}
g_{\epsilon}^{\mathrm{CK}}(X,Y)
=2K\operatorname{tr}(X^\top Y)-2\operatorname{tr}(X)\operatorname{tr}(Y)
+\epsilon\operatorname{tr}(X^\top Y),
\qquad \epsilon>0.
\label{eq:regularized_cartan_killing_preconditioner}
\end{equation}
The unregularized term is degenerate on the center of $\mathfrak{gl}(K)$, while the $\epsilon$ term makes the form positive definite. It is invariant under conjugation by $\mathrm{O}(K)$, not under general $\mathrm{GL}(K)$ conjugation, so it is an $\mathrm{O}(K)$-invariant conditioning choice rather than a full-$\mathrm{GL}(K)$-covariant natural gradient. The Fisher-Rao metric of the belief family and a chosen frame-coordinate metric are distinct geometric data. Equation~\eqref{eq:mu_gradient_general} already distinguishes the canonical and surrogate response terms through $\chi_{\mathrm{surr}}$; neither path contains an adaptive-$\alpha$ response term.

This bilinear form is distinct from the extrinsic ambient-Frobenius pullback through the exponential map,
\begin{equation}
g_{\phi}^{\mathrm{ext}}(X,Y)
=\operatorname{tr}\left(D\exp_{\phi}[X]^\top D\exp_{\phi}[Y]\right).
\label{eq:extrinsic_exp_pullback}
\end{equation}
The pullback is positive definite only where $D\exp_{\phi}$ has full rank. It is not generally left-invariant: for $h=\operatorname{diag}(2,1)$, left multiplication sends $E_{11}$ to $hE_{11}$ and changes its squared Frobenius norm from $1$ to $4$. A direct group-element alternative would be the left-invariant metric
\begin{equation}
g_U^{\mathrm{L}}(\xi,\eta)
=\operatorname{tr}\left[(U^{-1}\xi)^\top(U^{-1}\eta)\right],
\label{eq:left_invariant_group_metric}
\end{equation}
which is not bi-invariant and would require a group-element parameterization or an atlas. This alternative specifies geometry on group elements rather than on a fixed exponential chart.

If $\Sigma_i\approx\Omega_{ij}\Sigma_j\Omega_{ij}^\top$, Fisher preconditioning reduces the alignment mean term to $\beta_{ij}(\mu_i-\Omega_{ij}\mu_j)$. This covariance condition must be imposed or justified separately. Concentrated attention yields it from the stationary precision equation only when prior coupling is negligible or the prior covariance is aligned with the attended source. Agreement among transported neighbor precisions does not remove a mismatched prior.

\paragraph{Covariance gradient.}

The covariance channel has the same objective distinction. For an own-row filtering derivative, the target-blind result is

\begin{equation}
\boxed{
\begin{aligned}
g_{\Sigma_i}^{\mathrm{filt}}
&= \frac{1}{2}
\left[
-(1+\alpha_i^*) \Sigma_i^{-1}
+ \alpha_i^* \Sigma_{p,i}^{-1}
+ \sum_{j\in S_i} \beta_{ij}^*
(\Omega_{ij}\Sigma_j\Omega_{ij}^\top)^{-1}
\right] \\
&\quad+\chi_{\mathrm{surr}}
\sum_{j\in S_i}
\frac{\partial\beta_{ij}^*}{\partial\Sigma_i}E_{ij}.
\end{aligned}
}
\label{eq:Sigma_gradient_main}
\end{equation}

The precision aggregation is the direct derivative common to both objectives. The final response term belongs only to the entropy-suppressed surrogate. The optimized precision contributes $\alpha_i^*$ but no $\partial\alpha_i^*/\partial\Sigma_i$ term. A smoothing covariance derivative additionally includes all rows in which $\Sigma_i$ appears as a transported key. Setting the direct canonical own-row derivative to zero gives, for $\alpha_i^*=1$,

\begin{equation}
\Sigma_i^{-1}
=\frac{1}{2}
\left[
\Sigma_{p,i}^{-1}
+ \sum_{j\in S_i} \beta_{ij}^*
(\Omega_{ij}\Sigma_j\Omega_{ij}^\top)^{-1}
\right].
\label{eq:sigma_fixed_point_main}
\end{equation}

Each agent's precision is thus determined self-consistently as a
$\beta$-weighted combination of its prior precision and its neighbors'
transported precisions. When the state-dependent prior coupling is
weak ($\alpha_i^* \ll 1$; cf.\ Section~\ref{sec:state_dependent_precision}),
the prior term becomes negligible. Simultaneously, the entropy
contribution from the prior KL vanishes, changing the coefficient on
$\Sigma_i^{-1}$ from $-2$ to $-1$ in the gradient. The fixed point then
reduces to
$\Sigma_i^{-1} \approx
\langle(\Omega_{ij}\Sigma_j\Omega_{ij}^\top)^{-1}\rangle_\beta$,
the attention-weighted average of transported neighbor precisions.
More generally, with optimized state-dependent precision $\alpha_i^*$, the fixed point interpolates as
%
\begin{equation}
\Sigma_i^{-1}
= \frac{\alpha_i^*}{1+\alpha_i^*} \Sigma_{p,i}^{-1}
+ \frac{1}{1+\alpha_i^*}\sum_{j\in S_i} \beta_{ij}^*
(\Omega_{ij}\Sigma_j\Omega_{ij}^\top)^{-1},
\label{eq:sigma_fixed_point_alpha}
\end{equation}
%
a convex combination that continuously interpolates between
prior-dominated ($\alpha_i^* \gg 1$, precision $\approx \Sigma_{p,i}^{-1}$)
and alignment-dominated ($\alpha_i^* \ll 1$, precision
$\approx \langle(\Omega_{ij}\Sigma_j\Omega_{ij}^\top)^{-1}\rangle_\beta$)
regimes. Equation~\eqref{eq:state_dependent_alpha} makes $\alpha_i^*$ decrease as the full prior KL increases. This is the covariance counterpart of the adaptive anchoring described in Section~\ref{sec:state_dependent_precision}; it is not inferred from the singular mean-only limit.
The stationary precision is a weighted barycenter of the prior and transported neighbor precisions. Concentration of $\beta$ alone does not give pairwise covariance alignment. For one attended scalar source with variance $4$, prior variance $1$, and $\alpha=1$, the stationary variance is $2/(1+1/4)=8/5$, not $4$. Alignment requires negligible prior coupling or an aligned prior.

Section~\ref{app:covariance_dynamics} proves strict covariance convexity with both $\alpha$ and $\beta$ frozen. Optimizing either coordinate introduces a negative-semidefinite response term, so this conditional convexity does not prove local attractivity of the reduced dynamics. The covariance dynamics is an additional variational construction, separate from specifying a predictive mean operator.

\paragraph{Gauge frame gradient.}
The gauge frame gradient drives learning of the parallel transport operators $\Omega_{ij}=\exp(\phi_i)\exp(-\phi_j)$. Its structure involves the differential of the matrix exponential map and differs for compact ($\mathrm{SO}(N)$) and general ($\mathrm{GL}(K)$) gauge groups; Section~\ref{app:phi_gradients} derives both the canonical envelope form and the entropy-suppressed surrogate response.

\subsection{Identity-Plus-Correction Belief Updates}
\label{sec:residual_connections}

The variational update of a belief mean has an identity-plus-correction form. It must be distinguished from the explicit predictive-content skip connection in main-paper Eq.~\eqref{eq:relational_predictive_residual}. For either objective and either derivative-support mode, write the selected derivative as $g_i^{(\ell)}$. One step is

\begin{equation}
\mu_i^{(\ell+1)}
=\operatorname{Retr}_{\mu_i^{(\ell)}}
\left(-\eta\Sigma_i g_i^{(\ell)}\right),
\label{eq:residual_gradient_flow}
\end{equation}

with the ordinary step size $\eta$. For Euclidean mean coordinates the retraction is addition, so this has the form $\text{new}=\text{old}+\text{correction}$. Its resemblance to a transformer residual connection is structural; the variational construction does not uniquely derive the architecture's additive skip connection.

The optimized precision $\alpha_i^*$ controls the prior-restoring part of the correction, not the coefficient of the identity map itself. Large $\alpha_i^*$ strengthens anchoring to the prior; small $\alpha_i^*$ weakens that anchoring. This adaptive correction is not identical to a learned transformer skip gate.

For a belief coordinate $\theta_i\in\{\mu_i,\Sigma_i\}$, the common update notation is

\begin{equation}
\theta_i^{(\ell+1)}
=\operatorname{Retr}_{\theta_i^{(\ell)}}
\left(-\eta_\theta P_{\theta_i}g_{\theta_i}^{(\ell)}\right),
\end{equation}

where $P_{\theta_i}$ is the appropriate inverse-Fisher block and $g_{\theta_i}^{(\ell)}$ is the filtering or smoothing derivative of the selected objective. A frame update requires its own declared tangent metric; the regularized Cartan/Killing conditioner is not the Gaussian Fisher metric.

The alignment term favors agreement after transport: zero alignment cost requires $q_i=\Omega_{ij}q_j$, which becomes coordinate equality only in a common frame. Prior anchoring supplies a local restoring term that competes with that agreement when priors differ. Neither term alone guarantees a stable multi-cluster stationary state. The residual form above retains an identity path, while $\alpha_i$ changes the restoring correction rather than the coefficient of that identity map. These are algebraic properties of the update; long-time behavior requires a separate dynamical argument.

\paragraph{Layers and tied update-map iterations.}
A tied inner loop recomputes attention and applies the same belief-update map $T$ times. In smoothing mode, and under step-size and regularity conditions not established here, this can be interpreted as repeated descent on the selected global objective. In the optional own-row filtering mode it is only repeated application of the own-row update map; because key-side derivatives are excluded, it is not claimed to minimize a single global free energy or to decrease one monotonically. A finite composition of these maps is not an equilibrium solution without a specified fixed point and convergence or implicit-solution assumptions.


\section{Limits of Gradient-Field and Causal-Update Interpretations}
\label{sec:gradient_scope_limits}

An explicit predictive map need not be a variational descent field. The following examples separate a fixed-metric scalar-potential interpretation from an arbitrary coordinate transformation, and causal source support from the derivative support of a population objective.

\subsection{Coordinate MLPs and scalar-potential integrability}
\label{sec:mlp_gradient_scope}

An arbitrary coordinate MLP also need not be a natural-gradient field of a scalar objective. At fixed \(\Sigma\succ0\), a field \(f(\mu)=-\Sigma\nabla_\mu\Phi(\mu)\), with twice differentiable \(\Phi\), must satisfy symmetry of \(\Sigma^{-1}Df\). Consider a realizable single-direction field
\begin{equation}
f(\mu_1,\mu_2)=(0,a(\mu_1)),\qquad
\Sigma^{-1}=
\begin{pmatrix}u&v\\v&w\end{pmatrix}\succ0.
\label{eq:mlp_integrability_example}
\end{equation}
Its metric-weighted Jacobian is
\begin{equation}
\Sigma^{-1}Df=
\begin{pmatrix}
v a'(\mu_1)&0\\
w a'(\mu_1)&0
\end{pmatrix}.
\label{eq:mlp_integrability_jacobian}
\end{equation}
Since \(w>0\), symmetry fails wherever \(a'\ne0\), including for a GELU activation on such an interval. This excludes a universal fixed-metric scalar-potential interpretation of arbitrary MLP fields. It does not exclude specially constrained fields, augmented latent models, or a separately specified metric-dependent construction. The example concerns an instantaneous field, not every possible finite composite endpoint map.

\subsection{Causal source support and variational derivatives}
\label{sec:causal_gradient_scope}

A causal source graph alone does not make full population inference a causal update of every state. For the allowed edge \(2\leftarrow1\), take
\begin{equation}
E_{21}=\frac12(\mu_2-L_{21}\mu_1)^\top
R^{-1}(\mu_2-L_{21}\mu_1),\qquad R\succ0.
\end{equation}
At fixed link and covariance its source derivative is
\begin{equation}
\nabla_{\mu_1}E_{21}
=-L_{21}^{\top}R^{-1}(\mu_2-L_{21}\mu_1).
\label{eq:causal_source_derivative}
\end{equation}
If \(\mu_2\) has assimilated an observation unavailable at the earlier position, a full gradient update of \(\mu_1\) transmits that information backward. Omitting the source derivative defines a different directed update. This is the derivative-support distinction used in Section~\ref{app:optional_belief_revision}; a causal mask does not make the two vector fields equal.


The predictive alternative uses a state shared across candidate tokens and determined only by the available prefix, as specified in main-paper Appendix~\ref{app:architecture_scope}. Outer supervised differentiation updates shared parameters without making the next-token target an input to that prediction-time state.

% bibliographystyle is set by jmlr2e
\bibliography{../references}

\end{document}
